% !TeX program = pdflatex
% !BIB program = bibtex
% FisherBalancing v3: plain-language main text; explicit mT; proofs retained.
% Detailed proofs are included for author and independent mathematical review.
% This manuscript is not a formal proof certificate or a submission-format certification.
% Official template: https://aistats.org/aistats2027/AISTATS2027PaperPack.zip
\documentclass[10pt,twoside,letterpaper]{article}
% Use the official pack's fancyhdr when compiling from the project root.
% Register its standard package name before aistats2027 requires it.
\makeatletter
\IfFileExists{AISTATS/fancyhdr.sty}{%
  \usepackage{AISTATS/fancyhdr}%
  \@namedef{ver@fancyhdr.sty}{}%
}{}
\makeatother
\PassOptionsToPackage{round,authoryear}{natbib}
\usepackage{aistats2027}
\usepackage[T1]{fontenc}
\usepackage[utf8]{inputenc}
\usepackage{microtype}
\usepackage{amsmath,amssymb,mathtools,amsthm}
\usepackage{mathrsfs}
\usepackage{algorithm}
\usepackage{algorithmic}
\usepackage{graphicx}
\usepackage{subcaption}
\captionsetup[subfigure]{font=scriptsize,skip=1pt}
\usepackage{booktabs}
\usepackage{float}
\graphicspath{{AISTATS/figure_prepared/main/}{AISTATS/figure_prepared/appendix/}{figure_prepared/main/}{figure_prepared/appendix/}}
\usepackage{natbib}
\renewcommand{\bibsection}{\subsubsection*{\refname}}
\usepackage{xurl}
\usepackage{hyperref}
\hypersetup{hidelinks}
\usepackage[nameinlink,noabbrev]{cleveref}
\setlength{\emergencystretch}{1em}
\allowdisplaybreaks[1]

\newtheorem{theorem}{Theorem}
\newtheorem{proposition}{Proposition}
\newtheorem{lemma}{Lemma}[section]
\crefname{theorem}{Theorem}{Theorems}
\Crefname{theorem}{Theorem}{Theorems}
\crefname{proposition}{Proposition}{Propositions}
\Crefname{proposition}{Proposition}{Propositions}
\crefname{lemma}{Lemma}{Lemmas}
\Crefname{lemma}{Lemma}{Lemmas}
\newcommand{\R}{\mathbb R}
\newcommand{\E}{\mathbb E}
\newcommand{\Prob}{\mathbb P}
\newcommand{\tr}{\operatorname{tr}}
\newcommand{\Cov}{\operatorname{Cov}}
\newcommand{\op}{\mathrm{op}}
\newcommand{\argmin}{\operatorname*{arg\,min}}
\newcommand{\argmax}{\operatorname*{arg\,max}}
\newcommand{\norm}[1]{\lVert #1\rVert}
\newcommand{\FB}{\textnormal{\textsc{FisherBalancing}}}
\newcommand{\thetahat}{\widehat\theta}
\newcommand{\Vhat}{\widehat V}
\newcommand{\gstar}{g^\star}
\newcommand{\Rdir}{\mathfrak R_{\mathrm{dir},T}}
\newcommand{\Rtwo}{\mathfrak R_{2,T}}
\newcommand{\PaperTitle}{FisherBalancing: Learning the Optimal Information Rate\protect\\ in Batched GLM Estimation}

\begin{document}
\runningtitle{Learning the Optimal Information Rate in Batched GLM Estimation}
\runningauthor{Anonymous Author}
\twocolumn[
\aistatstitle{\PaperTitle}
\aistatsauthor{Anonymous Author}
\aistatsaddress{Anonymous Institution}
]

\begin{abstract}
To accurately estimate parameters with a limited number of observations,
it is important to decide which data to observe. Information matrices are
often used for this purpose, but when the information matrix depends on
the unknown true parameter, a good parameter estimate is needed first in
order to select good data. We study this problem in Generalized Linear
Models (GLMs) and propose \FB{}, which uses the current estimate to
continually reinforce directions with less information. Even when
selecting multiple data points at once, \FB{} avoids repeatedly selecting
similar information and collects sufficient information across all
parameter directions without requiring an accurate initial estimate or
separate forced exploration. As a result, information in the
least-informed direction grows at the best achievable rate, and we
further provide guarantees on finite-time parameter error and the optimal
asymptotic error in the hardest direction. Experiments using synthetic
GLMs and inputs obtained from real datasets evaluate both information
balance and parameter-estimation error.
\end{abstract}


\section{Introduction}
\label{sec:introduction}

With a limited observation budget, parameter estimation requires deciding
\emph{which inputs to select for observing responses}. Optimal experimental
design, adaptive estimation, active regression, and active learning address
this question \cite{Pukelsheim2006,HuanJagalurMarzouk2024}. Each selected
input $x$ is a \emph{query}; observing its response $y$ yields a data point
$(x,y)$. The \emph{information matrix} measures how much these data tell us
about different parameter directions. Under standard conditions, its
inverse describes estimation uncertainty. Classical A-, D-, and
E-optimality use different properties of this matrix to guide selection.

A natural adaptive approach estimates the parameter from collected data
and uses the resulting information matrix to select new queries. Such methods
are studied in sequential and adaptive experimental design \cite{Freise2021,Freise2024}. The difficulty is that
the information matrix can depend on the unknown true parameter, even in
generalized linear models (GLMs) \cite{KhuriEtAl2006,HuanJagalurMarzouk2024}.
This creates a circular dependence: selecting useful queries requires a
good parameter estimate, but obtaining that estimate requires useful
queries. We study this approach for GLMs and ask:

\begin{quote}
\emph{Starting from an inaccurate parameter estimate, can we keep using
the current estimate to select queries and still collect the information needed
for accurate estimation sufficiently quickly?}
\end{quote}

Accurate estimation of the full parameter requires more than a large
amount of information on average. A poorly informed direction---a
particular combination of parameters---can remain highly uncertain. We
therefore aim to strengthen the least-informed directions. This principle
corresponds to \emph{E-optimality}, which maximizes the smallest eigenvalue
of the information matrix. Whereas D-optimality measures overall
information and A-optimality measures average estimation uncertainty,
E-optimality directly targets the direction with the least information.

We propose \FB{} to implement this principle. Using the information matrix
computed from the current estimate, it scores queries by how much
information they add to poorly informed directions. It weights multiple
directions, giving more weight to those with less information, rather than
following only the weakest one. Within a batch, it adds each selected
query's expected information to the matrix before choosing the next query.
This reduces redundant selections without observing new responses. On a
finite candidate pool, the method simply alternates between scoring
candidates and updating the matrix.

An inaccurate estimate can misjudge which queries are useful, so the
success of this rule is not immediate. Under our assumptions, we show
that information accumulates in every direction without accurate
initialization or queries reserved solely for exploration. The parameter
estimate and the resulting information matrix then become more accurate.
Eventually, information in the least-informed direction grows at the
best achievable rate.

We also connect this information growth to estimation error. We bound
parameter error with finitely many observations and, under additional
conditions, show that the largest mean squared error over unit parameter
directions attains the optimal asymptotic constant.

Our contributions are threefold.
\begin{itemize}
    \item \textbf{Optimal information growth from an inaccurate estimate.}
    Without accurate initialization or separate forced exploration,
    information accumulates in every direction and eventually grows at
    the optimal rate in the least-informed direction.

    \item \textbf{A simple adaptive E-optimal rule.}
    \FB{} weights poorly informed directions and accounts for the
    information expected from queries already selected within a batch.

    \item \textbf{Guarantees for parameter estimation.}
    We bound finite-sample parameter error and, under additional
    conditions, establish the optimal asymptotic constant for
    worst-direction mean squared error.
\end{itemize}

Experiments on synthetic GLMs and inputs from real datasets compare
\FB{} with A-, D-, and E-optimal design, bandit-based methods, and
active-learning methods. We evaluate both information in the weakest
direction and overall parameter-estimation error.










\section{Related Work}
\label{sec:related}

\paragraph{Optimal experimental design.}
Optimal experimental design allocates a limited observation budget to
estimate parameters accurately. A-, D-, and E-optimality use different
properties of the information matrix \cite{Pukelsheim2006}. Recent work
covers nonlinear and non-Gaussian models, sequential design, and
computational methods \cite{HuanJagalurMarzouk2024}, including continuous
relaxation, rounding, exchange methods, and regret minimization
\cite{AllenZhu2017}.
Our focus is adaptive GLM design, where queries are selected using an
information matrix computed from the current parameter estimate.

We specifically target E-optimality, which maximizes the smallest
information eigenvalue. Selecting a fixed-size subset of candidates is
combinatorial, and some variants are NP-hard
\cite{BrownLaddhaSingh2024}. Optimizing sampling proportions instead
allows convex or semidefinite formulations. Rather than solving a
subset-selection problem, \FB{} repeatedly reinforces poorly informed
directions. In our GLM setting, where queries can be repeated, this rule
eventually achieves optimal information growth in the weakest direction.

\paragraph{Adaptive estimation and regression.}
Adaptive experimental design uses the current parameter estimate to
choose the next query. Adaptive Wynn-type methods analyze this
process in nonlinear regression and GLMs, establishing consistency and
long-run D-optimality \cite{Freise2021,Freise2024}. Active regression
instead selects which responses to obtain from available inputs when
labels are expensive \cite{ChenPrice2019,ChenLiSun2022}. Its goal is to
approach the prediction or regression accuracy of the full dataset with
fewer labels.

\FB{} targets E-optimal information growth for the full GLM parameter.
Under our assumptions, information accumulates in every direction even
from an inaccurate estimate, and growth in the weakest direction
eventually becomes optimal. For batches, the same rule selects queries
one at a time rather than solving a separate joint design problem.

\paragraph{Active learning and preference learning.}
Active learning selects which inputs to label using uncertainty, diversity,
or model gradients. BADGE combines uncertainty and diversity
\cite{Ash2020}; BAIT uses Fisher information and its connection to
estimation error for batch selection \cite{Ash2021}; and FIRAL connects
Fisher-based selection with prediction error in logistic classification
\cite{ChenBiros2023}. While many such methods target prediction or
classification, \FB{} targets balanced estimation of the full GLM
parameter, with guarantees on weakest-direction information growth and
parameter error.

Preference learning offers a natural application. In a pairwise logistic
model, the feature difference between two alternatives forms a GLM query,
so choosing a comparison determines which parameter directions are
learned. Recent work selects informative comparisons for direct preference
optimization (DPO), reinforcement learning from human feedback (RLHF),
and reward modeling \cite{Muldrew2024,Das2025,Shen2025}.
\FB{} can prioritize comparisons that inform poorly known parameter
directions, supporting a preference model intended for use across
multiple decisions.
% ============================================================================
% ============================================================================
% ============================================================================
% ============================================================================
\section{Problem Setup}
\label{sec:setup}

We first define the batched GLM estimation problem, the Fisher information
used to evaluate queries, and the three quantities that we will track
throughout the paper.

\subsection{Queries and batched responses}

We begin with the interaction protocol: the learner repeatedly chooses
queries, observes their responses, and updates an unknown GLM parameter.

Let $\theta^\star\in\Theta\subset\R^d$ be the unknown parameter and let
$\mathcal X\subset\R^d$ be the feasible query set. We keep
$\mathcal X$ general here and specialize its geometry only in
Section~\ref{sec:theory}.

Fix a batch size $m\ge1$. At round $t$, the learner uses only previous
responses to choose
\begin{equation}
\mathcal B_t
=
(x_{t,1},\ldots,x_{t,m})
\in\mathcal X^m,
\end{equation}
and then observes all $m$ responses together. Thus $m=1$ gives
sequential feedback, while $m>1$ gives batched feedback. After $T$
rounds, the total number of responses is $mT$.

Responses follow a regular canonical generalized linear model (GLM):
\begin{equation}
p_\theta(y\mid x)
=
h(y)\exp\!\left\{
yx^\top\theta-A(x^\top\theta)
\right\}.
\label{eq:glm}
\end{equation}
Its conditional mean is
\begin{equation}
\mu_\theta(x)=A'(x^\top\theta).
\label{eq:glm-mean}
\end{equation}
Examples include Bernoulli-logistic and fixed-variance Gaussian
regression.

Let $\mathcal D_t$ contain all observations from the first $t$ rounds.
For ridge coefficient $\rho>0$, we compute
\begin{equation}
\begin{split}
\thetahat_t
\in
\argmin_{\theta\in\Theta}
\Bigg\{&
\sum_{s=1}^{t}\sum_{j=1}^{m}
\left[
A(x_{s,j}^\top\theta)
-
Y_{s,j}x_{s,j}^\top\theta
\right]\\
&+
\frac{\rho}{2}\norm{\theta}_2^2
\Bigg\}.
\end{split}
\label{eq:mle}
\end{equation}
Round $t$ uses $\thetahat_{t-1}$ to select its batch and updates
$\thetahat_t$ only after the full batch is observed.

Thus the learner must choose useful queries using only its current
parameter estimate, with no feedback available inside a batch.


\subsection{Fisher information}

We next quantify what a query teaches us about the parameter. Fisher
information provides this quantity directly from the GLM likelihood.

Differentiating the GLM log-likelihood gives the score
\begin{equation}
g_\theta(x,y)
=
\nabla_\theta\log p_\theta(y\mid x)
=
\bigl(y-\mu_\theta(x)\bigr)x.
\label{eq:likelihood-score}
\end{equation}
Since
$\operatorname{Var}_\theta(Y\mid x)=A''(x^\top\theta)$,
the Fisher information from query $x$ is
\begin{equation}
\begin{aligned}
G_\theta(x)
&=
\E_\theta[
g_\theta(x,Y)g_\theta(x,Y)^\top
\mid x]\\
&=
w_\theta(x)xx^\top,
\qquad
w_\theta(x):=A''(x^\top\theta).
\end{aligned}
\label{eq:fisher}
\end{equation}
For a batch $\mathcal B=(x_1,\ldots,x_m)$,
\begin{equation}
G_\theta(\mathcal B)
=
\sum_{j=1}^{m}G_\theta(x_j).
\label{eq:batch-fisher}
\end{equation}

For a unit direction $u$,
\begin{equation}
u^\top G_\theta(x)u
=
w_\theta(x)(x^\top u)^2.
\label{eq:directional-information}
\end{equation}
Thus an informative query must both align with the direction and have
sufficient response curvature.

The selected queries define the true and plug-in information matrices
\begin{equation}
\begin{aligned}
V_t^\star
&=
\rho I+
\sum_{s=1}^{t}
G_{\theta^\star}(\mathcal B_s),\\
\Vhat_t
&=
\rho I+
\sum_{s=1}^{t}
G_{\thetahat_t}(\mathcal B_s).
\end{aligned}
\label{eq:states}
\end{equation}
Both use the same queries. $V_t^\star$ is evaluated at the unknown true
parameter and is used only for analysis, whereas $\Vhat_t$ is
computable from the current estimate.

The algorithm will therefore select queries using $\Vhat_t$, while our
analysis checks whether this computable information also reflects the
true information $V_t^\star$.


\subsection{Performance criteria}

We now specify what ``good estimation'' means. We track both how quickly
information grows and how this growth translates into parameter error.

For any positive definite matrix $V$,
\begin{equation}
\max_{\norm{u}_2=1}
u^\top V^{-1}u
=
\lambda_{\min}(V)^{-1}.
\label{eq:inverse-radius}
\end{equation}
Thus $\lambda_{\min}(V)$ measures information in the least-covered
parameter direction.

Let $\mathcal P(\mathcal X)$ be the set of probability measures on
$\mathcal X$. For $\nu\in\mathcal P(\mathcal X)$, define
\begin{equation}
M(\nu)
=
\int_{\mathcal X}
G_{\theta^\star}(x)\,d\nu(x).
\end{equation}
The optimal information growth rate per query is
\begin{equation}
\begin{aligned}
\gamma^\star
&=
\sup_{\nu\in\mathcal P(\mathcal X)}
\lambda_{\min}(M(\nu)),\\
\gstar
&=
m\gamma^\star.
\end{aligned}
\label{eq:gamma-star}
\end{equation}
Hence $\gamma^\star$ is the best long-run rate at which any query
allocation can increase its weakest directional information.

Let
\begin{equation}
e_T
=
\thetahat_T-\theta^\star.
\end{equation}
We evaluate the Euclidean and worst-direction RMSEs
\begin{equation}
\begin{aligned}
\Rtwo
&=
\left(
\E_{\theta^\star}
\norm{e_T}_2^2
\right)^{1/2},\\
\Rdir
&=
\sup_{\norm{u}_2=1}
\left(
\E_{\theta^\star}
[(u^\top e_T)^2]
\right)^{1/2}.
\end{aligned}
\label{eq:risks}
\end{equation}
If
$Q_T=\E_{\theta^\star}[e_Te_T^\top]$,
then
\begin{equation}
\Rtwo^2
=
\tr(Q_T),
\qquad
\Rdir^2
=
\lambda_{\max}(Q_T).
\end{equation}

Our main performance quantities are
$\gamma^\star$, $\Rdir$, and $\Rtwo$: information growth, error in the
hardest direction, and total parameter error.


% ============================================================================
\section{Algorithm}
\label{sec:algorithm}

We now introduce \FB{}. The key idea is to give more attention to
directions with less information and to update this balance while a batch
is being constructed.

\subsection{Weighting under-observed directions}

We first define how \FB{} identifies which parameter directions currently
need more information.

Recall that $\Vhat_t$ is the plug-in information matrix. While building
a new batch, we maintain a temporary plug-in state $W$, initialized as
$W=\Vhat_{t-1}$.

For temperature $\tau>0$, define
\begin{equation}
P_\tau(W)
=
\frac{\exp(-W/\tau)}
{\tr\exp(-W/\tau)}.
\label{eq:gibbs}
\end{equation}
If $\lambda_i$ is an eigenvalue of $W$, its eigenvector receives weight
proportional to $e^{-\lambda_i/\tau}$. Thus less-informed directions
receive larger weight.

For an available estimate $\widehat\theta$, define
\begin{equation}
\begin{aligned}
s(x;W,\widehat\theta)
&=
\tr\!\left(
P_\tau(W)G_{\widehat\theta}(x)
\right)\\
&=
w_{\widehat\theta}(x)
x^\top P_\tau(W)x.
\end{aligned}
\label{eq:score}
\end{equation}
The score measures the information expected to be added to currently
under-observed directions. Equivalently, for any positive semidefinite
trace-one matrix $P$,
\begin{equation}
\E_\theta
\left[
\norm{P^{1/2}g_\theta(x,Y)}_2^2
\mid x
\right]
=
\tr(PG_\theta(x)).
\label{eq:projected-score}
\end{equation}

Thus each candidate receives a large score when it is informative
precisely where the current estimate has accumulated less information.


\subsection{Greedy batch construction}

We next use this score to construct a batch without observing any
responses inside the batch.

At round $t$, set $W=\Vhat_{t-1}$. For $j=1,\ldots,m$, choose
\begin{equation}
x_{t,j}
\in
\argmax_{x\in\mathcal X}
\tr\!\left(
P_{\tau_t}(W)
G_{\thetahat_{t-1}}(x)
\right),
\label{eq:batch-query}
\end{equation}
then update
\begin{equation}
W
\leftarrow
W+
G_{\thetahat_{t-1}}(x_{t,j}).
\label{eq:temporary-update}
\end{equation}
This update uses no new response; it only accounts for information
expected from queries already selected in the current batch. After all
$m$ queries are chosen, their responses are observed together and the
model is refitted.

\begin{algorithm}[t]
\caption{\FB{} with batch size $m$}
\label{alg:fb}
\begin{algorithmic}[1]
\REQUIRE Query set $\mathcal X$; parameter set $\Theta$; batch size $m$;
ridge $\rho>0$; fixed coefficient $\tau_0>0$
(prescribed-budget calibration: \eqref{eq:temperature}).
\STATE Set $\mathcal D_0=\varnothing$, $\Vhat_0=\rho I$, and
$\thetahat_0\in\argmin_{\theta\in\Theta}\norm{\theta}_2^2$.
\FOR{$t=1,\ldots,T$}
  \STATE Set $\tau_t=\tau_0(mt)^{3/4}$ and $W=\Vhat_{t-1}$.
  \FOR{$j=1,\ldots,m$}
    \STATE Select $x_{t,j}$ by \eqref{eq:batch-query}.
    \STATE Update $W$ by \eqref{eq:temporary-update}.
  \ENDFOR
  \STATE Observe $Y_{t,1},\ldots,Y_{t,m}$ together.
  \STATE Refit $\thetahat_t$ and recompute $\Vhat_t$.
\ENDFOR
\RETURN $\thetahat_T$ and $\Vhat_T$.
\end{algorithmic}
\end{algorithm}

The temporary update is what makes the method naturally batched: each
new query is chosen after accounting for what the rest of the current
batch is already expected to learn.


\subsection{Connection to optimal information growth}

Finally, we explain why the acquisition score is aligned with our target
$\gamma^\star$, rather than being only a heuristic weighting rule.

The target $\gamma^\star$ is the best achievable growth rate of the
minimum directional information. To connect our score to this objective,
define
\begin{equation}
\Phi_\tau(W)
=
-\tau\log\tr\exp(-W/\tau).
\label{eq:potential}
\end{equation}
It satisfies
\begin{equation}
\lambda_{\min}(W)-\tau\log d
\le
\Phi_\tau(W)
\le
\lambda_{\min}(W),
\end{equation}
and
\begin{equation}
\left.
\frac{d}{d\eta}
\Phi_\tau
\left(
W+\eta G_\theta(x)
\right)
\right|_{\eta=0}
=
\tr\!\left(
P_\tau(W)G_\theta(x)
\right).
\label{eq:potential-gradient}
\end{equation}
Thus the \FB{} score is the first-order improvement of a smooth
minimum-information objective.

For a prescribed calibration budget $N_{\max}=mT_{\max}$ and a
positive reference scale $R$, we set, for $d\ge2$,
\begin{equation}
\begin{aligned}
\tau_{\mathrm{range}}
&=R\sqrt{\frac{m}{8\log d}},\\
\tau_0
&=\frac{\tau_{\mathrm{range}}}{N_{\max}^{3/4}},
\qquad
\tau_t=\tau_0(mt)^{3/4}.
\end{aligned}
\label{eq:temperature}
\end{equation}
The temperature reaches $\tau_{\mathrm{range}}$ at the prescribed
budget. Appendix~\ref{app:range-calibration} gives the derivation and
its applicability conditions. The asymptotic results below hold
$\tau_0$ fixed.

In short, \FB{} greedily improves a smooth version of the same
least-direction information criterion used to define
$\gamma^\star$.


% ============================================================================
\section{Theory}
\label{sec:theory}

We now ask three questions: how fast information can possibly grow,
whether \FB{} reaches that rate, and whether this information growth is
statistically optimal for parameter estimation.

The algorithm is defined for general $\mathcal X$. For analysis, we now
specialize to the full Euclidean ball
\begin{equation}
\mathcal X
=
\left\{
x\in\R^d:
\norm{x}_2\le L
\right\}.
\label{eq:theory-ball}
\end{equation}
Queries may be repeated with fresh responses. The following guarantees
apply to this full-ball setting; finite candidate pools used in the
experiments are empirical extensions.


\subsection{Conditions}
\label{sec:conditions}

We first state the regularity conditions needed to control both the GLM
estimator and the information computed from its current estimate.

Assume $\Theta$ is compact and convex, contains $\theta^\star$, and
\begin{equation}
\sup_{\theta\in\Theta}
\norm{\theta}_2
\le
S.
\end{equation}
For all $z,z'\in[-LS,LS]$, assume
\begin{equation}
\begin{aligned}
0<\underline w
&\le
A''(z)
\le
\overline w<\infty,\\
|A''(z)-A''(z')|
&\le
L_A|z-z'|.
\end{aligned}
\label{eq:regularity}
\end{equation}

With
\begin{equation}
\xi_{t,j}
=
Y_{t,j}
-
\mu_{\theta^\star}(x_{t,j}),
\end{equation}
assume conditional independence within each selected batch and
\begin{equation}
\E
\left[
e^{a\xi_{t,j}}
\mid
\mathcal D_{t-1},\mathcal B_t
\right]
\le
e^{a^2\sigma^2/2},
\qquad
a\in\R.
\label{eq:noise-condition}
\end{equation}

These assumptions cover bounded logistic and unit-variance Gaussian
GLMs. We further assume exact maximization of
\eqref{eq:batch-query} over the ball. We do not require orthogonal
queries, an eigenvalue gap, accurate initialization, or forced
exploration.

These conditions are sufficient to learn the unknown GLM while keeping
the plug-in Fisher information uniformly controlled.


\subsection{Optimal information growth}
\label{sec:information-rate}

We first identify the fastest rate any policy could achieve, and then
compare \FB{} against this universal benchmark.

Under the full-ball setting, the supremum defining $\gamma^\star$ is
attained.

\begin{proposition}[Universal information bound]
\label{prop:ceiling}
For every adaptive policy,
\begin{equation}
\frac{\underline wL^2}{d}
\le
\gamma^\star
\le
\frac{\overline wL^2}{d},
\end{equation}
and
\begin{equation}
\lambda_{\min}(V_T^\star)
\le
\rho+mT\gamma^\star.
\label{eq:ceiling}
\end{equation}
\end{proposition}

Thus $\gamma^\star$ is a universal upper rate for minimum directional
information.

Define
\begin{equation}
\varepsilon_{m,T}(\delta)
=
C
\left[
\frac{\log(emT/\delta)}
{(mT)^{1/4}}
+
\frac1T
\right],
\label{eq:rate-error}
\end{equation}
where $C$ is independent of $m,T,\delta$.

\begin{theorem}[\FB{} information optimality]
\label{thm:information}
With probability at least $1-\delta$, simultaneously for all $T\ge1$,
\begin{equation}
\left|
\frac{\lambda_{\min}(\Vhat_T)}{mT}
-
\gamma^\star
\right|
\le
\varepsilon_{m,T}(\delta),
\label{eq:plugin-rate}
\end{equation}
and
\begin{equation}
0
\le
\frac{
\rho+mT\gamma^\star
-
\lambda_{\min}(V_T^\star)
}{
mT
}
\le
\varepsilon_{m,T}(\delta).
\label{eq:true-rate}
\end{equation}
\end{theorem}

Thus both the plug-in and true information approach the optimal rate.
Moreover, once
$\varepsilon_{m,T}(\delta)\le\gamma^\star/2$,
\begin{equation}
\frac{
\norm{\Vhat_T-V_T^\star}_{\op}
}{
mT
}
\le
C'
\sqrt{
\frac{
\log(emT/\delta)
}{
mT
}
}.
\label{eq:tracking}
\end{equation}

For fixed $m$,
\begin{equation}
\begin{aligned}
\frac{\lambda_{\min}(\Vhat_T)}{T}
&\longrightarrow
\gstar,\\
\frac{\lambda_{\min}(V_T^\star)}{T}
&\longrightarrow
\gstar
\end{aligned}
\qquad
\text{a.s.}
\label{eq:main-limit}
\end{equation}

The $1/T$ term reflects fewer parameter updates under larger batches;
it is an analysis upper bound, not an unavoidable batching cost. We do
not claim that the finite-time rate $(mT)^{-1/4}$ is optimal.

The main conclusion is that \FB{} asymptotically collects information
in its hardest direction as fast as any policy can.


\subsection{Parameter recovery}
\label{sec:confidence}

We next translate the information guarantee into a direct bound on the
error of the estimated parameter.

Define
\begin{equation}
C_{\mathrm{glm}}
=
\frac{
\sqrt{\max\{1,\overline w\}}
}{
\min\{1,\underline w\}
},
\end{equation}
and
\begin{equation}
\begin{split}
\beta_{m,T}(\delta)
=
C_{\mathrm{glm}}
\Bigg[
&
\sigma
\sqrt{
d\log\!\left(
1+\frac{mTL^2}{\rho d}
\right)
+
2\log\frac2\delta
}\\
&+
\sqrt\rho\,S
\Bigg].
\end{split}
\label{eq:confidence-radius}
\end{equation}

\begin{proposition}[\FB{} parameter recovery]
\label{prop:confidence}
For every unit $u$,
\begin{equation}
|u^\top e_T|
\le
\beta_{m,T}(\delta)
\sqrt{
u^\top\Vhat_T^{-1}u
}.
\label{eq:directional-confidence}
\end{equation}
If $\varepsilon_{m,T}(\delta)<\gamma^\star$, then
\begin{equation}
\norm{e_T}_2
\le
\frac{
\beta_{m,T}(\delta)
}{
\sqrt{
mT
\left(
\gamma^\star
-
\varepsilon_{m,T}(\delta)
\right)
}
}.
\label{eq:l2-confidence}
\end{equation}
\end{proposition}

Both bounds hold on the event of
Theorem~\ref{thm:information}. For fixed $m$,
$\thetahat_T\to\theta^\star$ almost surely.

Thus the optimal information growth above is not only a spectral
property: it directly yields shrinking directional and Euclidean
parameter error.


\subsection{Sharp worst-direction risk}
\label{sec:sharp-risk}

Finally, we ask whether the worst-direction error of \FB{} is only
consistent, or whether its asymptotic constant is also the best possible.

We first give a lower bound for competing adaptive estimators, and then
show that \FB{} with the ridge MLE attains it.

For policy $\pi$, define
\begin{equation}
\mathcal J_T^\pi
=
\E_{\theta^\star}
\left[
\sum_{t=1}^{T}
G_{\theta^\star}(\mathcal B_t^\pi)
\right].
\end{equation}
For a competing estimator $\widetilde\theta_T$, let
\begin{equation}
a_T(\theta)
=
\E_\theta[\widetilde\theta_T],
\qquad
D_T
=
\nabla_\theta a_T(\theta^\star).
\end{equation}

We consider estimators with finite second moments for which $a_T$ is
differentiable at an interior $\theta^\star$ and likelihood
differentiation is valid. Define
\begin{equation}
U_T
=
\sum_{t=1}^{T}\sum_{j=1}^{m}
x_{t,j}
\left[
Y_{t,j}
-
\mu_{\theta^\star}(x_{t,j})
\right].
\end{equation}
The required identity is
\begin{equation}
D_T
=
\E_{\theta^\star}
\left[
(\widetilde\theta_T-\theta^\star)
U_T^\top
\right].
\label{eq:mean-score-identity}
\end{equation}

\begin{proposition}[Universal directional lower bound]
\label{prop:risk-lower}
Every estimator in this comparison class satisfies
\begin{equation}
\mathfrak R_{\mathrm{dir},T}
(\widetilde\theta_T)^2
\ge
\frac{
s_{\min}(D_T)^2
}{
mT\gamma^\star
}.
\label{eq:risk-lower-general}
\end{equation}
\end{proposition}

In particular, local unbiasedness gives
\begin{equation}
\mathfrak R_{\mathrm{dir},T}^2
\ge
\frac1{mT\gamma^\star},
\end{equation}
while $D_T\to I$ implies
\begin{equation}
\liminf_{T\to\infty}
T\gstar\,
\mathfrak R_{\mathrm{dir},T}^2
\ge
1.
\label{eq:asymptotic-lower}
\end{equation}

This gives the benchmark: no asymptotically regular estimator in the
comparison class can beat the constant $1$ after normalization by
$T\gstar$.

We now return to \FB{}. Assume additionally that $\theta^\star$ is
interior to $\Theta$ and that every measure attaining $\gamma^\star$
induces the same information matrix. For any maximizing measure
$\nu^\star$, write
\begin{equation}
M^\star
=
M(\nu^\star),
\qquad
J^\star
=
mM^\star.
\end{equation}

\begin{proposition}[\FB{} attains the sharp bound]
\label{prop:efficiency}
For \FB{} with the ridge MLE and fixed $m$,
\begin{equation}
mTQ_T
\longrightarrow
(M^\star)^{-1},
\qquad
D_T
\longrightarrow
I.
\label{eq:mse-matrix}
\end{equation}
\end{proposition}

Consequently,
\begin{equation}
T\gstar\,\Rdir^2
\longrightarrow
1,
\label{eq:sharp-directional-risk}
\end{equation}
and
\begin{equation}
T\Rtwo^2
\longrightarrow
\tr\!\left(
(J^\star)^{-1}
\right).
\label{eq:sharp-euclidean-risk}
\end{equation}
Together with Proposition~\ref{prop:risk-lower}, this shows that
\FB{} attains the optimal leading constant for worst-direction mean
squared error within the stated comparison class.

Finally,
\begin{equation}
\frac1{\gstar}
\le
\tr\!\left(
(J^\star)^{-1}
\right)
\le
\frac d{\gstar}.
\end{equation}
We claim sharp optimality for worst-direction risk, not A-optimality of
the total Euclidean risk.

Thus the theory progresses from a universal information benchmark to
finite-sample recovery and finally to sharp asymptotic optimality in
the hardest parameter direction.

All proofs are given in Appendix~\ref{app:proofs}.

\section{Experiments}
\label{sec:experiments}

We compare \textbf{Ours} with A-, D-, and E-optimal design, UCB,
Thompson sampling, and uniform sampling. Figures report minimum true
information and Euclidean parameter error against $mt$ additional
queries after a 40-query D-optimal initialization; $t$ counts batch
rounds. Curves are seed means with pointwise $\pm1.96\,\mathrm{SEM}$
bands. Appendix~\ref{app:experiment-protocol} details the implementations.

\subsection{Synthetic Experiment}
\label{sec:synthetic-experiment}

Anisotropic unit-ball pools vary in dimension, covariance decay,
target orientation, and batch size. Each condition uses ten paired seeds.

\begin{figure}[!htbp]
  \centering
  \includegraphics[width=\columnwidth]{synthetic_gaussian_bernoulli.pdf}
  \caption{Synthetic pools with $d=20$, $\alpha=1.5$, a weak target,
  and $m=10$. Columns compare Gaussian and Bernoulli responses;
  rows show minimum information (higher is better) and parameter
  error (lower is better). Bands summarize ten paired seeds.}
  \label{fig:synthetic-main}
\end{figure}

\Cref{fig:synthetic-main} compares matched Gaussian and Bernoulli
responses. At $mt=200$, Ours reduces error by $6.5\%$ and $6.0\%$
relative to the best baseline in each family. Minimum information is
nearly tied with the best Gaussian baseline and improves by $2.0\%$
for Bernoulli responses. Appendix~\ref{app:synthetic-details} gives the
selection rule; Appendix~\ref{app:synthetic-results} reports all displayed sweeps.

\subsection{Empirical Experiment}
\label{sec:empirical-experiment}

We use eight OpenML covariate pools, three response families, and
$m\in\{1,5,10,20\}$. Original labels are replaced by generated GLM
responses; each condition uses three paired run seeds.

\begin{figure}[!htbp]
  \centering
  \includegraphics[width=\columnwidth]{empirical_gaussian_poisson.pdf}
  \caption{Australian covariates with $d=20$ and $m=5$.
  Columns compare Gaussian and Poisson responses on the same pool
  and target; rows show minimum information and parameter error,
  as in \cref{fig:synthetic-main}. Bands summarize three paired run
  seeds. Appendix~\ref{app:empirical-results} reports the broader grid.}
  \label{fig:empirical-main}
\end{figure}

In \cref{fig:empirical-main}, Ours reduces Gaussian and Poisson error
by $12.1\%$ and $4.0\%$ at $mt=200$. Information improves by $0.9\%$
for Gaussian responses and is $0.9\%$ below the best baseline for Poisson
responses. Each percentage compares with the best baseline on that metric.

Across all 96 conditions, Ours leads on information in 48, error in 33,
and both in 24. Both main figures show favorably selected cases;
the broader appendix results use the same seven algorithms.

\clearpage
\subsection*{AI Use Statement}
% TODO (authors): Extend this statement to cover any earlier AI assistance
% with research, proofs, code, experiments, or manuscript writing.
Generative AI tools were used to adapt the manuscript to the AISTATS 2027
\LaTeX{} template, consolidate the bibliography, inspect existing
experiment outputs, prepare plotting code and figures, and draft the
experimental descriptions.

\bibliographystyle{plainnat}
\IfFileExists{AISTATS/reference_final.bib}
  {\bibliography{AISTATS/reference_final}}
  {\bibliography{reference_final}}

\section*{Checklist}
% TODO (authors): Select Yes, No, or Not Applicable for each official question
% and add a section reference or brief justification where appropriate.

\begin{enumerate}

  \item For all models and algorithms presented, check if you include:
  \begin{enumerate}
    \item A clear description of the mathematical setting, assumptions, algorithm, and/or model. [Yes/No/Not Applicable]
    \item An analysis of the properties and complexity (time, space, sample size) of any algorithm. [Yes/No/Not Applicable]
    \item (Optional) Anonymized source code, with specification of all dependencies, including external libraries. [Yes/No/Not Applicable]
  \end{enumerate}

  \item For any theoretical claim, check if you include:
  \begin{enumerate}
    \item Statements of the full set of assumptions of all theoretical results. [Yes/No/Not Applicable]
    \item Complete proofs of all theoretical results. [Yes/No/Not Applicable]
    \item Clear explanations of any assumptions. [Yes/No/Not Applicable]
  \end{enumerate}

  \item For all figures and tables that present empirical results, check if you include:
  \begin{enumerate}
    \item The code, data, and instructions needed to reproduce the main experimental results (either in the supplemental material or as a URL). [Yes/No/Not Applicable]
    \item All the training details (e.g., data splits, hyperparameters, how they were chosen). [Yes/No/Not Applicable]
    \item A clear definition of the specific measure or statistics and error bars (e.g., with respect to the random seed after running experiments multiple times). [Yes/No/Not Applicable]
    \item A description of the computing infrastructure used. (e.g., type of GPUs, internal cluster, or cloud provider). [Yes/No/Not Applicable]
  \end{enumerate}

  \item If you are using existing assets (e.g., code, data, models) or curating/releasing new assets, check if you include:
  \begin{enumerate}
    \item Citations of the creator If your work uses existing assets. [Yes/No/Not Applicable]
    \item The license information of the assets, if applicable. [Yes/No/Not Applicable]
    \item New assets either in the supplemental material or as a URL, if applicable. [Yes/No/Not Applicable]
    \item Information about consent from data providers/curators. [Yes/No/Not Applicable]
    \item Discussion of sensible content if applicable, e.g., personally identifiable information or offensive content. [Yes/No/Not Applicable]
  \end{enumerate}

  \item If you used crowdsourcing or conducted research with human subjects, check if you include:
  \begin{enumerate}
    \item The full text of instructions given to participants and screenshots. [Yes/No/Not Applicable]
    \item Descriptions of potential participant risks, with links to Institutional Review Board (IRB) approvals if applicable. [Yes/No/Not Applicable]
    \item The estimated hourly wage paid to participants and the total amount spent on participant compensation. [Yes/No/Not Applicable]
  \end{enumerate}

\end{enumerate}

\clearpage
\appendix
\onecolumn
\thispagestyle{empty}
\aistatstitle{\PaperTitle\\Supplementary Materials}
\numberwithin{equation}{section}
\section{Additional Experiments and Details}
\label{app:experiments}

This appendix describes the finite-pool implementations used in
\cref{sec:experiments}, states how the main panels were selected, and
reports the prepared synthetic sweeps and empirical trajectory grids.
The empirical main panels reappear in the Australian grids, and the
synthetic main panels in the batch-size sweeps, alongside other conditions.

\subsection{Evaluation Protocol and Implementation}
\label{app:experiment-protocol}

\paragraph{Metrics, initialization, and budget.}
Every method starts with the same label-free, greedy D-optimal selection
of $n_0=40$ inputs, using unit curvature for initialization. Responses
are then generated and a ridge estimate is fitted with $\rho=1$.
We retain the manuscript's batch-round index $t$: $mt$ is the number
of additional queries and $n_0+mt$ is the total response count.
The initial information includes the ridge term and the warm-start
queries $x_{0,i}$. At each completed batch boundary, we record
\begin{equation}
  \begin{aligned}
    V_0^\star&=I+\sum_{i=1}^{n_0}
      w_{\theta^\star}(x_{0,i})x_{0,i}x_{0,i}^\top,\qquad
    V_t^\star=V_0^\star+\sum_{s=1}^t G_{\theta^\star}(\mathcal B_s),\\
    I_t&=\lambda_{\min}(V_t^\star),\qquad
    e_t=\norm{\thetahat_t-\theta^\star}_2.
  \end{aligned}
  \label{eq:experimental-metrics}
\end{equation}
Here $\thetahat_t$ is fitted to all $n_0+mt$ observed responses.
The true parameter is used to evaluate $I_t$; acquisition decisions
use the current fitted parameter. Runs continue to 400 additional
responses. All figures display the same prefix $0\le mt\le200$.
Values are recorded after $t=0,1,2,\ldots$ completed rounds and plotted
at $mt=0,m,2m,\ldots$; lines connect these measurements.

\paragraph{Implementation of Ours.}
The archived results reported in this section use the implementation
below, not the prescribed-budget calibration in \eqref{eq:temperature}.
The label Ours and its blue line are fixed across all response families
and batch sizes. For $m=1$, the implementation selects the largest value of
$w_{\widehat\theta}(x)(x^\top v_{\min})^2$, where $v_{\min}$ is one
minimum-eigenvalue eigenvector of the plug-in information matrix.
For $m>1$, it uses the Gibbs score in \eqref{eq:score}. Its fitted parameter
is frozen during a batch, and its temporary information matrix is
updated after each selected input. The archived runs use
\begin{equation}
  \tau=\max\!\left\{10^{-8},\ R\sqrt{\frac{m}{8\log d}}\right\}.
  \label{eq:experimental-temperature}
\end{equation}
For Gaussian and Bernoulli responses on the unit-radius pools,
$R=1$ and $R=1/4$, respectively. For Poisson responses, $R$ is the
largest candidate value of
$w_{\widehat\theta}(x)\norm{x}_2^2$, recomputed over the remaining
candidates within a batch. The temperature multiplier is one.

Inputs cannot be repeated within a batch.
Synthetic pools permit reuse across batches; empirical
pools are sampled without replacement.

\paragraph{Relation to the theoretical setting.}
These runs use finite candidate pools, range-based soft temperatures,
and unconstrained ridge fits. The theoretical analysis uses maximization
over the full ball, the growing temperature in \cref{alg:fb}, and a
compact parameter constraint. Poisson responses also extend beyond
the global sub-Gaussian noise assumption. The experiments therefore
assess the information-balancing mechanism under the stated practical
choices; their endpoint improvements are finite-budget empirical results.

\paragraph{Optimization and comparators.}
Gaussian ridge regression is solved directly. Logistic and Poisson
fits use damped Newton updates with at most 80 iterations and tolerance
$10^{-8}$. Acquisition curvature clips linear predictors to $[-30,30]$
and uses floors of $10^{-10}$ for logistic and $10^{-12}$ for Poisson
curvature. Poisson fitting also guards the exponent at 30. The
generated Poisson predictors lie within $[-8.5,7.5]$, so exponent
clipping is inactive for the generated responses and true-information
metric. Acquisition considers at most 2,000 available candidates per
batch, selected using a method-independent priority for shared remaining
rows. Common comparators are greedy A-, D-, and E-optimal design,
UCB/GLM-UCB, Thompson sampling, and uniform random sampling. E-optimal
design evaluates the one-step minimum eigenvalue on all candidates when
there are at most 400; otherwise it uses the top 256 hard-score candidates
plus 32 random candidates. The same six baselines and Ours appear in
every panel, for all response families and batch sizes. All methods
use the same feedback protocol and initialization within a condition.

\paragraph{Seeds and uncertainty.}
Synthetic conditions use seeds $0,\ldots,9$; empirical conditions use
seeds $0,1,2$. Curves are pointwise arithmetic means and shading is
mean $\pm1.96\,\mathrm{SEM}$. These are descriptive uncertainty bands,
with only three run seeds in the empirical setting. Response
randomness is paired across methods by observation index for synthetic
pools and by row identity for empirical pools. Within each empirical
dataset, the feature preprocessing and target are fixed across runs.
Run seeds also control initialization and acquisition randomization;
these bands do not measure uncertainty over datasets.

\subsection{Empirical Pools and Main-Panel Selection}
\label{app:empirical-selection}

\paragraph{Real covariates with generated responses.}
The eight processed OpenML pools are listed in
\cref{tab:empirical-universe}. Original labels are discarded. Numeric
columns undergo median imputation and standardization; categorical
columns use most-frequent imputation and bounded one-hot encoding.
PCA for dense matrices or truncated SVD for sparse matrices reduces
the feature representation, after which reduced coordinates are
standardized and every row is projected to the unit ball. The
preprocessing seed is zero. The table reports the actual processed
row counts and dimensions in the completed-run metadata.

The target follows a minimum-eigenvalue eigenvector of the empirical
second-moment matrix $X^\top X/|X|$ and is scaled toward unit predictor standard deviation,
with Euclidean norm capped at eight. Responses have unit-variance
Gaussian noise, the Bernoulli-logistic link, or Poisson mean
$\exp(-0.5+x^\top\theta^\star)$. Each dataset is evaluated under all
three families and $m\in\{1,5,10,20\}$, yielding 96 conditions.

\begin{table}[htbp]
\centering
\begin{tabular}{lrrrrrr}
\toprule
Dataset & OpenML ID & $d$ & Rows & Info. & Error & Both\\
\midrule
Australian      & 40981 & 20 &   690 &  8 & 8 & 5\\
PhishingWebsites &  4534 & 20 & 10000 & 10 & 2 & 2\\
adult           &  1590 & 20 & 10000 &  6 & 5 & 4\\
bank-marketing  &  1461 & 20 & 10000 &  5 & 3 & 2\\
credit-g        &    31 & 20 &  1000 &  9 & 5 & 5\\
eeg-eye-state   &  1471 & 12 & 10000 &  1 & 0 & 0\\
spambase        &    44 & 20 &  4601 &  9 & 7 & 6\\
splice          &    46 & 20 &  3190 &  0 & 3 & 0\\
\midrule
All settings    &       &    &       & 48 & 33 & 24\\
\bottomrule
\end{tabular}
\caption{Empirical selection universe. Each dataset contributes 12
conditions. The last three columns count strict improvements in
three-seed endpoint means over the strongest of the six common
comparators at $mt=200$, for information, error, and both metrics.
Ours uses the implementation specified for each batch size above.
All panels use this same comparison set.}
\label{tab:empirical-universe}
\end{table}

\paragraph{Selection rule.}
We pair a Gaussian condition with a Bernoulli or Poisson condition on
the same dataset at the same batch size $m\in\{5,10,20\}$, giving 48
candidate pairs. Let $\overline I_{f,\mathrm{Ours}}$ and
$\overline e_{f,\mathrm{Ours}}$ denote mean endpoints at $mt=200$ in
family $f$, and let $\mathcal C$ contain the six common design,
bandit, and random comparators. We retain pairs whose error mean is
smaller than every member of $\mathcal C$ in both families, then maximize
\begin{equation}
  S_{\mathrm{pair}}=\sqrt{S_{\mathrm{Gaussian}}S_f},\qquad
  S_f=\sqrt{
    \frac{\overline I_{f,\mathrm{Ours}}}
         {\max_{c\in\mathcal C}\overline I_{f,c}}
    \frac{\min_{c\in\mathcal C}\overline e_{f,c}}
         {\overline e_{f,\mathrm{Ours}}}}.
  \label{eq:empirical-selection}
\end{equation}
This selects Australian with Gaussian and Poisson responses at $m=5$.
Information dominance is not required by this rule. All runs enter
the means, and the main figure displays Ours and the same six common
comparators in both families. The appendix uses the identical algorithm
set and reports both information and parameter error.

For Gaussian responses, Ours has
$(\overline I,\overline e)=(12.0704,1.1090)$; the strongest common
information and error comparators are A-optimal design ($11.9585$)
and D-optimal design ($1.2616$), respectively. For Poisson responses,
Ours gives $(4.6865,1.7021)$; the corresponding comparators are
E-optimal design ($4.7313$) and A-optimal design ($1.7728$).

Across all 96 conditions, the median relative information gain over the
best common comparator is approximately zero, and median relative
error reduction is $-1.33\%$. The favorable main panels therefore
illustrate conditions where information balancing helps; they do not
summarize typical gains over the entire empirical grid.
The following trajectory pages retain Australian and credit-g from
the prepared report and show every response family and batch size for
those two datasets. \Cref{tab:empirical-universe} covers all eight pools.

\subsection{Synthetic Construction and Sweeps}
\label{app:synthetic-details}

For each seed we draw 3,000 anisotropic Gaussian candidates with
coordinate standard deviations $j^{-\alpha/2}$. Eight percent are
replaced by directions supported on the final four coordinates. Rows
are projected to the Euclidean unit ball, and the $2d$ signed coordinate
anchors are appended. The resulting pool has $3000+2d$ inputs.

Let $u_{\min}$ and $u_{\max}$ be eigenvectors of the empirical second-moment
matrix $X^\top X/|X|$ for its smallest and largest eigenvalues.
Weak and strong targets use these directions;
a random target uses an independent normalized Gaussian vector. For
each direction $u$, we set
\begin{equation}
  \theta^\star=cu,\qquad
  c=\min\{\operatorname{std}(Xu)^{-1},8\}.
  \label{eq:experimental-target}
\end{equation}
The pool is regenerated across synthetic seeds, so the eigenvector-aligned
targets also vary across seeds. Gaussian responses have noise standard
deviation one; Bernoulli responses use the logistic link. The additional
Poisson comparison uses the offset specified above.

The central condition is $d=20$, $\alpha=1.5$, and a weak target.
The prepared Gaussian and logistic sweeps vary one factor at a time:
target orientation (weak, random, strong), dimension ($10,20,40$),
covariance decay ($0.5,1.5,3$), and batch size ($1,5,10,20$).
The batch-size figures display all four values $m=1,5,10,20$.
Batch-size panels hold the central geometry fixed; additional batched
geometry runs remain in the experiment archive.
The Poisson comparison contains the central condition with $m=1$ and
$m=10$. All figures compare the same seven algorithms.

\paragraph{Synthetic main-panel selection.}
The displayed slices contain 12 matched pairs: ten Gaussian/Bernoulli
pairs and two Gaussian/Poisson pairs. Each pair fixes the geometry,
target construction, and batch size. We retain pairs in which Ours
improves both endpoint metrics over every baseline in both families,
then maximize the smaller of the two relative error reductions.
All ten seeds enter the means. This selects Gaussian and Bernoulli
responses in the central setting with $m=10$ for
\cref{fig:synthetic-main}. At $mt=200$, information gains are $0.05\%$
and $1.95\%$, and error reductions are $6.50\%$ and $5.98\%$,
respectively. This is a favorable illustrative pair selected from the
displayed slices; the appendix also includes conditions with smaller
or negative gains.

\paragraph{Reproducing the displayed figures.}
The experiment archive retains each run's configuration, completion
marker, metadata, and trajectory. The accompanying
\texttt{scripts/prepare\_aistats\_figures.py} regenerates these panels
from the completed trajectories. Its figure manifest records source
files, hashes, seed and checkpoint checks, selected conditions, and
endpoint summaries. The 96-condition comparison uses 2,016 completed
trajectories (seven methods including the evaluated variant, with
three seeds per condition).

\clearpage
\subsection{Synthetic Results}
\label{app:synthetic-results}

% Each source PDF is one metric panel with its own internal legend.
% Change these subfigure widths or row breaks to rearrange the LaTeX grids.
\newcommand{\experimentpanel}[3][0.485]{%
  \begin{subfigure}[t]{#1\linewidth}%
    \centering
    \includegraphics[width=\linewidth]{#2}%
    \caption{#3}%
  \end{subfigure}%
}

\begin{figure}[H]
  \centering
  \experimentpanel{synthetic_gaussian_d20_alpha1.5_weak_m1_lambda_min.pdf}{Weak target: minimum information.}\hfill
  \experimentpanel{synthetic_gaussian_d20_alpha1.5_weak_m1_l2_error.pdf}{Weak target: parameter error.}
  \par\vspace{2pt}
  \experimentpanel{synthetic_gaussian_d20_alpha1.5_random_m1_lambda_min.pdf}{Random target: minimum information.}\hfill
  \experimentpanel{synthetic_gaussian_d20_alpha1.5_random_m1_l2_error.pdf}{Random target: parameter error.}
  \par\vspace{2pt}
  \experimentpanel{synthetic_gaussian_d20_alpha1.5_strong_m1_lambda_min.pdf}{Strong target: minimum information.}\hfill
  \experimentpanel{synthetic_gaussian_d20_alpha1.5_strong_m1_l2_error.pdf}{Strong target: parameter error.}
  \caption{Gaussian target-orientation sweep at $d=20$, $\alpha=1.5$, and $m=1$. Each row changes the target direction; columns show information and parameter error. Curves summarize ten seeds.}
  \label{fig:synthetic-theta-gaussian}
\end{figure}
\clearpage
\begin{figure}[H]
  \centering
  \experimentpanel{synthetic_logistic_d20_alpha1.5_weak_m1_lambda_min.pdf}{Weak target: minimum information.}\hfill
  \experimentpanel{synthetic_logistic_d20_alpha1.5_weak_m1_l2_error.pdf}{Weak target: parameter error.}
  \par\vspace{2pt}
  \experimentpanel{synthetic_logistic_d20_alpha1.5_random_m1_lambda_min.pdf}{Random target: minimum information.}\hfill
  \experimentpanel{synthetic_logistic_d20_alpha1.5_random_m1_l2_error.pdf}{Random target: parameter error.}
  \par\vspace{2pt}
  \experimentpanel{synthetic_logistic_d20_alpha1.5_strong_m1_lambda_min.pdf}{Strong target: minimum information.}\hfill
  \experimentpanel{synthetic_logistic_d20_alpha1.5_strong_m1_l2_error.pdf}{Strong target: parameter error.}
  \caption{Bernoulli-logistic target-orientation sweep at $d=20$, $\alpha=1.5$, and $m=1$. Each row changes the target direction; columns show information and parameter error. Curves summarize ten seeds.}
  \label{fig:synthetic-theta-logistic}
\end{figure}
\clearpage
\begin{figure}[H]
  \centering
  \experimentpanel{synthetic_gaussian_d10_alpha1.5_weak_m1_lambda_min.pdf}{$d=10$: minimum information.}\hfill
  \experimentpanel{synthetic_gaussian_d10_alpha1.5_weak_m1_l2_error.pdf}{$d=10$: parameter error.}
  \par\vspace{2pt}
  \experimentpanel{synthetic_gaussian_d20_alpha1.5_weak_m1_lambda_min.pdf}{$d=20$: minimum information.}\hfill
  \experimentpanel{synthetic_gaussian_d20_alpha1.5_weak_m1_l2_error.pdf}{$d=20$: parameter error.}
  \par\vspace{2pt}
  \experimentpanel{synthetic_gaussian_d40_alpha1.5_weak_m1_lambda_min.pdf}{$d=40$: minimum information.}\hfill
  \experimentpanel{synthetic_gaussian_d40_alpha1.5_weak_m1_l2_error.pdf}{$d=40$: parameter error.}
  \caption{Gaussian dimension sweep, with covariance decay $\alpha=1.5$, weak targets, and $m=1$. Information and error are compared within each dimension using ten seeds.}
  \label{fig:synthetic-dimension-gaussian}
\end{figure}
\clearpage
\begin{figure}[H]
  \centering
  \experimentpanel{synthetic_logistic_d10_alpha1.5_weak_m1_lambda_min.pdf}{$d=10$: minimum information.}\hfill
  \experimentpanel{synthetic_logistic_d10_alpha1.5_weak_m1_l2_error.pdf}{$d=10$: parameter error.}
  \par\vspace{2pt}
  \experimentpanel{synthetic_logistic_d20_alpha1.5_weak_m1_lambda_min.pdf}{$d=20$: minimum information.}\hfill
  \experimentpanel{synthetic_logistic_d20_alpha1.5_weak_m1_l2_error.pdf}{$d=20$: parameter error.}
  \par\vspace{2pt}
  \experimentpanel{synthetic_logistic_d40_alpha1.5_weak_m1_lambda_min.pdf}{$d=40$: minimum information.}\hfill
  \experimentpanel{synthetic_logistic_d40_alpha1.5_weak_m1_l2_error.pdf}{$d=40$: parameter error.}
  \caption{Bernoulli-logistic dimension sweep, with covariance decay $\alpha=1.5$, weak targets, and $m=1$. Information and error are compared within each dimension using ten seeds.}
  \label{fig:synthetic-dimension-logistic}
\end{figure}
\clearpage
\begin{figure}[H]
  \centering
  \experimentpanel{synthetic_gaussian_d20_alpha0.5_weak_m1_lambda_min.pdf}{$\alpha=0.5$: minimum information.}\hfill
  \experimentpanel{synthetic_gaussian_d20_alpha0.5_weak_m1_l2_error.pdf}{$\alpha=0.5$: parameter error.}
  \par\vspace{2pt}
  \experimentpanel{synthetic_gaussian_d20_alpha1.5_weak_m1_lambda_min.pdf}{$\alpha=1.5$: minimum information.}\hfill
  \experimentpanel{synthetic_gaussian_d20_alpha1.5_weak_m1_l2_error.pdf}{$\alpha=1.5$: parameter error.}
  \par\vspace{2pt}
  \experimentpanel{synthetic_gaussian_d20_alpha3_weak_m1_lambda_min.pdf}{$\alpha=3$: minimum information.}\hfill
  \experimentpanel{synthetic_gaussian_d20_alpha3_weak_m1_l2_error.pdf}{$\alpha=3$: parameter error.}
  \caption{Gaussian covariance-decay sweep at $d=20$ with weak targets and $m=1$. Larger $\alpha$ gives more unequal coordinate scales before projection. Each condition uses ten seeds.}
  \label{fig:synthetic-covariance-gaussian}
\end{figure}
\clearpage
\begin{figure}[H]
  \centering
  \experimentpanel{synthetic_logistic_d20_alpha0.5_weak_m1_lambda_min.pdf}{$\alpha=0.5$: minimum information.}\hfill
  \experimentpanel{synthetic_logistic_d20_alpha0.5_weak_m1_l2_error.pdf}{$\alpha=0.5$: parameter error.}
  \par\vspace{2pt}
  \experimentpanel{synthetic_logistic_d20_alpha1.5_weak_m1_lambda_min.pdf}{$\alpha=1.5$: minimum information.}\hfill
  \experimentpanel{synthetic_logistic_d20_alpha1.5_weak_m1_l2_error.pdf}{$\alpha=1.5$: parameter error.}
  \par\vspace{2pt}
  \experimentpanel{synthetic_logistic_d20_alpha3_weak_m1_lambda_min.pdf}{$\alpha=3$: minimum information.}\hfill
  \experimentpanel{synthetic_logistic_d20_alpha3_weak_m1_l2_error.pdf}{$\alpha=3$: parameter error.}
  \caption{Bernoulli-logistic covariance-decay sweep at $d=20$ with weak targets and $m=1$. Larger $\alpha$ gives more unequal coordinate scales before projection. Each condition uses ten seeds.}
  \label{fig:synthetic-covariance-logistic}
\end{figure}
\clearpage
\begin{figure}[H]
  \centering
  \experimentpanel[0.465]{synthetic_gaussian_d20_alpha1.5_weak_m1_lambda_min.pdf}{$m=1$: minimum information.}\hfill
  \experimentpanel[0.465]{synthetic_gaussian_d20_alpha1.5_weak_m1_l2_error.pdf}{$m=1$: parameter error.}
  \par\vspace{2pt}
  \experimentpanel[0.465]{synthetic_gaussian_d20_alpha1.5_weak_m5_lambda_min.pdf}{$m=5$: minimum information.}\hfill
  \experimentpanel[0.465]{synthetic_gaussian_d20_alpha1.5_weak_m5_l2_error.pdf}{$m=5$: parameter error.}
  \par\vspace{2pt}
  \experimentpanel[0.465]{synthetic_gaussian_d20_alpha1.5_weak_m10_lambda_min.pdf}{$m=10$: minimum information.}\hfill
  \experimentpanel[0.465]{synthetic_gaussian_d20_alpha1.5_weak_m10_l2_error.pdf}{$m=10$: parameter error.}
  \par\vspace{2pt}
  \experimentpanel[0.465]{synthetic_gaussian_d20_alpha1.5_weak_m20_lambda_min.pdf}{$m=20$: minimum information.}\hfill
  \experimentpanel[0.465]{synthetic_gaussian_d20_alpha1.5_weak_m20_l2_error.pdf}{$m=20$: parameter error.}
  \caption{Gaussian batch-size sweep at the central geometry. All panels use the same query budget $mt$, and each curve summarizes ten seeds. Ours has the same label and color at every batch size; measurements follow completed batches. The $m=10$ row reproduces the corresponding family in \cref{fig:synthetic-main}.}
  \label{fig:synthetic-batch-gaussian}
\end{figure}
\clearpage
\begin{figure}[H]
  \centering
  \experimentpanel[0.465]{synthetic_logistic_d20_alpha1.5_weak_m1_lambda_min.pdf}{$m=1$: minimum information.}\hfill
  \experimentpanel[0.465]{synthetic_logistic_d20_alpha1.5_weak_m1_l2_error.pdf}{$m=1$: parameter error.}
  \par\vspace{2pt}
  \experimentpanel[0.465]{synthetic_logistic_d20_alpha1.5_weak_m5_lambda_min.pdf}{$m=5$: minimum information.}\hfill
  \experimentpanel[0.465]{synthetic_logistic_d20_alpha1.5_weak_m5_l2_error.pdf}{$m=5$: parameter error.}
  \par\vspace{2pt}
  \experimentpanel[0.465]{synthetic_logistic_d20_alpha1.5_weak_m10_lambda_min.pdf}{$m=10$: minimum information.}\hfill
  \experimentpanel[0.465]{synthetic_logistic_d20_alpha1.5_weak_m10_l2_error.pdf}{$m=10$: parameter error.}
  \par\vspace{2pt}
  \experimentpanel[0.465]{synthetic_logistic_d20_alpha1.5_weak_m20_lambda_min.pdf}{$m=20$: minimum information.}\hfill
  \experimentpanel[0.465]{synthetic_logistic_d20_alpha1.5_weak_m20_l2_error.pdf}{$m=20$: parameter error.}
  \caption{Bernoulli-logistic batch-size sweep at the central geometry. All panels use the same query budget $mt$, and each curve summarizes ten seeds. Ours has the same label and color at every batch size; measurements follow completed batches. The $m=10$ row reproduces the corresponding family in \cref{fig:synthetic-main}.}
  \label{fig:synthetic-batch-logistic}
\end{figure}
\clearpage
\begin{figure}[H]
  \centering
  \experimentpanel{synthetic_poisson_m1_lambda_min.pdf}{$m=1$: minimum information.}\hfill
  \experimentpanel{synthetic_poisson_m1_l2_error.pdf}{$m=1$: parameter error.}
  \par\vspace{2pt}
  \experimentpanel{synthetic_poisson_m10_lambda_min.pdf}{$m=10$: minimum information.}\hfill
  \experimentpanel{synthetic_poisson_m10_l2_error.pdf}{$m=10$: parameter error.}
  \caption{Poisson responses in the central synthetic setting, with $m=1$ and $m=10$. The batched implementation uses the temperature in \eqref{eq:experimental-temperature}. Each condition uses ten seeds.}
  \label{fig:synthetic-poisson}
\end{figure}
\clearpage
\subsection{Empirical Results}
\label{app:empirical-results}

\begin{figure}[H]
  \centering
  \experimentpanel{empirical_Australian_gaussian_m1_lambda_min.pdf}{Gaussian: minimum information.}\hfill
  \experimentpanel{empirical_Australian_gaussian_m1_l2_error.pdf}{Gaussian: parameter error.}
  \par\vspace{2pt}
  \experimentpanel{empirical_Australian_logistic_m1_lambda_min.pdf}{Bernoulli: minimum information.}\hfill
  \experimentpanel{empirical_Australian_logistic_m1_l2_error.pdf}{Bernoulli: parameter error.}
  \par\vspace{2pt}
  \experimentpanel{empirical_Australian_poisson_m1_lambda_min.pdf}{Poisson: minimum information.}\hfill
  \experimentpanel{empirical_Australian_poisson_m1_l2_error.pdf}{Poisson: parameter error.}
  \caption{Australian covariates with $m=1$. Rows use Gaussian, Bernoulli, and Poisson responses on the same feature pool and weak target; columns show minimum information and parameter error. Curves summarize three paired run seeds.}
  \label{fig:empirical-australian-1}
\end{figure}
\clearpage
\begin{figure}[H]
  \centering
  \experimentpanel{empirical_Australian_gaussian_m5_lambda_min.pdf}{Gaussian: minimum information.}\hfill
  \experimentpanel{empirical_Australian_gaussian_m5_l2_error.pdf}{Gaussian: parameter error.}
  \par\vspace{2pt}
  \experimentpanel{empirical_Australian_logistic_m5_lambda_min.pdf}{Bernoulli: minimum information.}\hfill
  \experimentpanel{empirical_Australian_logistic_m5_l2_error.pdf}{Bernoulli: parameter error.}
  \par\vspace{2pt}
  \experimentpanel{empirical_Australian_poisson_m5_lambda_min.pdf}{Poisson: minimum information.}\hfill
  \experimentpanel{empirical_Australian_poisson_m5_l2_error.pdf}{Poisson: parameter error.}
  \caption{Australian covariates with $m=5$. Rows use Gaussian, Bernoulli, and Poisson responses on the same feature pool and weak target; columns show minimum information and parameter error. Curves summarize three paired run seeds. The Gaussian and Poisson rows reproduce the four panels of \cref{fig:empirical-main}.}
  \label{fig:empirical-australian-5}
\end{figure}
\clearpage
\begin{figure}[H]
  \centering
  \experimentpanel{empirical_Australian_gaussian_m10_lambda_min.pdf}{Gaussian: minimum information.}\hfill
  \experimentpanel{empirical_Australian_gaussian_m10_l2_error.pdf}{Gaussian: parameter error.}
  \par\vspace{2pt}
  \experimentpanel{empirical_Australian_logistic_m10_lambda_min.pdf}{Bernoulli: minimum information.}\hfill
  \experimentpanel{empirical_Australian_logistic_m10_l2_error.pdf}{Bernoulli: parameter error.}
  \par\vspace{2pt}
  \experimentpanel{empirical_Australian_poisson_m10_lambda_min.pdf}{Poisson: minimum information.}\hfill
  \experimentpanel{empirical_Australian_poisson_m10_l2_error.pdf}{Poisson: parameter error.}
  \caption{Australian covariates with $m=10$. Rows use Gaussian, Bernoulli, and Poisson responses on the same feature pool and weak target; columns show minimum information and parameter error. Curves summarize three paired run seeds.}
  \label{fig:empirical-australian-10}
\end{figure}
\clearpage
\begin{figure}[H]
  \centering
  \experimentpanel{empirical_Australian_gaussian_m20_lambda_min.pdf}{Gaussian: minimum information.}\hfill
  \experimentpanel{empirical_Australian_gaussian_m20_l2_error.pdf}{Gaussian: parameter error.}
  \par\vspace{2pt}
  \experimentpanel{empirical_Australian_logistic_m20_lambda_min.pdf}{Bernoulli: minimum information.}\hfill
  \experimentpanel{empirical_Australian_logistic_m20_l2_error.pdf}{Bernoulli: parameter error.}
  \par\vspace{2pt}
  \experimentpanel{empirical_Australian_poisson_m20_lambda_min.pdf}{Poisson: minimum information.}\hfill
  \experimentpanel{empirical_Australian_poisson_m20_l2_error.pdf}{Poisson: parameter error.}
  \caption{Australian covariates with $m=20$. Rows use Gaussian, Bernoulli, and Poisson responses on the same feature pool and weak target; columns show minimum information and parameter error. Curves summarize three paired run seeds.}
  \label{fig:empirical-australian-20}
\end{figure}
\clearpage
\begin{figure}[H]
  \centering
  \experimentpanel{empirical_credit-g_gaussian_m1_lambda_min.pdf}{Gaussian: minimum information.}\hfill
  \experimentpanel{empirical_credit-g_gaussian_m1_l2_error.pdf}{Gaussian: parameter error.}
  \par\vspace{2pt}
  \experimentpanel{empirical_credit-g_logistic_m1_lambda_min.pdf}{Bernoulli: minimum information.}\hfill
  \experimentpanel{empirical_credit-g_logistic_m1_l2_error.pdf}{Bernoulli: parameter error.}
  \par\vspace{2pt}
  \experimentpanel{empirical_credit-g_poisson_m1_lambda_min.pdf}{Poisson: minimum information.}\hfill
  \experimentpanel{empirical_credit-g_poisson_m1_l2_error.pdf}{Poisson: parameter error.}
  \caption{credit-g covariates with $m=1$. Rows use Gaussian, Bernoulli, and Poisson responses on the same feature pool and weak target; columns show minimum information and parameter error. Curves summarize three paired run seeds.}
  \label{fig:empirical-creditg-1}
\end{figure}
\clearpage
\begin{figure}[H]
  \centering
  \experimentpanel{empirical_credit-g_gaussian_m5_lambda_min.pdf}{Gaussian: minimum information.}\hfill
  \experimentpanel{empirical_credit-g_gaussian_m5_l2_error.pdf}{Gaussian: parameter error.}
  \par\vspace{2pt}
  \experimentpanel{empirical_credit-g_logistic_m5_lambda_min.pdf}{Bernoulli: minimum information.}\hfill
  \experimentpanel{empirical_credit-g_logistic_m5_l2_error.pdf}{Bernoulli: parameter error.}
  \par\vspace{2pt}
  \experimentpanel{empirical_credit-g_poisson_m5_lambda_min.pdf}{Poisson: minimum information.}\hfill
  \experimentpanel{empirical_credit-g_poisson_m5_l2_error.pdf}{Poisson: parameter error.}
  \caption{credit-g covariates with $m=5$. Rows use Gaussian, Bernoulli, and Poisson responses on the same feature pool and weak target; columns show minimum information and parameter error. Curves summarize three paired run seeds.}
  \label{fig:empirical-creditg-5}
\end{figure}
\clearpage
\begin{figure}[H]
  \centering
  \experimentpanel{empirical_credit-g_gaussian_m10_lambda_min.pdf}{Gaussian: minimum information.}\hfill
  \experimentpanel{empirical_credit-g_gaussian_m10_l2_error.pdf}{Gaussian: parameter error.}
  \par\vspace{2pt}
  \experimentpanel{empirical_credit-g_logistic_m10_lambda_min.pdf}{Bernoulli: minimum information.}\hfill
  \experimentpanel{empirical_credit-g_logistic_m10_l2_error.pdf}{Bernoulli: parameter error.}
  \par\vspace{2pt}
  \experimentpanel{empirical_credit-g_poisson_m10_lambda_min.pdf}{Poisson: minimum information.}\hfill
  \experimentpanel{empirical_credit-g_poisson_m10_l2_error.pdf}{Poisson: parameter error.}
  \caption{credit-g covariates with $m=10$. Rows use Gaussian, Bernoulli, and Poisson responses on the same feature pool and weak target; columns show minimum information and parameter error. Curves summarize three paired run seeds.}
  \label{fig:empirical-creditg-10}
\end{figure}
\clearpage
\begin{figure}[H]
  \centering
  \experimentpanel{empirical_credit-g_gaussian_m20_lambda_min.pdf}{Gaussian: minimum information.}\hfill
  \experimentpanel{empirical_credit-g_gaussian_m20_l2_error.pdf}{Gaussian: parameter error.}
  \par\vspace{2pt}
  \experimentpanel{empirical_credit-g_logistic_m20_lambda_min.pdf}{Bernoulli: minimum information.}\hfill
  \experimentpanel{empirical_credit-g_logistic_m20_l2_error.pdf}{Bernoulli: parameter error.}
  \par\vspace{2pt}
  \experimentpanel{empirical_credit-g_poisson_m20_lambda_min.pdf}{Poisson: minimum information.}\hfill
  \experimentpanel{empirical_credit-g_poisson_m20_l2_error.pdf}{Poisson: parameter error.}
  \caption{credit-g covariates with $m=20$. Rows use Gaussian, Bernoulli, and Poisson responses on the same feature pool and weak target; columns show minimum information and parameter error. Curves summarize three paired run seeds.}
  \label{fig:empirical-creditg-20}
\end{figure}
\clearpage

\section{Theoretical Proofs and Details}
\label{app:proofs}

The proof order is non-circular: deterministic excitation is established
for arbitrary fitted parameters, then concentration gives preliminary
localization, and only then is optimal information growth proved.
Mean-square efficiency is proved last. Throughout this appendix,
$N=mT$ denotes a completed response count; a general prefix length
$n$ need not be a multiple of $m$.
For completed counts $N=mT$, we write $\varepsilon_{N,m}(\delta)$
and $\beta_N(\delta)$ for $\varepsilon_{m,T}(\delta)$ and
$\beta_{m,T}(\delta)$ in \eqref{eq:rate-error} and~\eqref{eq:confidence-radius},
respectively. This changes notation only, not either bound.

\subsection{Details of Notation and Adaptive Scores}
\label{app:notation}

\subsubsection*{Indices, matrices, and constants.}
The flat index $i=(t-1)m+j$ identifies query $j$ in batch $t$; set
$x_i=x_{t,j}$, $Y_i=Y_{t,j}$, and
$\xi_i=Y_i-\mu_{\theta^\star}(x_i)$. Define the regularized input
matrix, the true information, and the accumulated likelihood score by
\begin{equation}
H_n=\rho I+\sum_{i=1}^n x_ix_i^\top,\qquad
Z_n=\rho I+\sum_{i=1}^nG_{\theta^\star}(x_i),\qquad
U_n=\sum_{i=1}^n x_i\xi_i.
\label{eq:app-states}
\end{equation}
Empty sums are zero, so $H_0=Z_0=\rho I$ and $U_0=0$.
At a completed round, $Z_{mt}=V_t^\star$. For a positive definite
matrix $H$, write $\norm{v}_H=(v^\top Hv)^{1/2}$.
The notation $A_1\preceq A_2$ means that $A_2-A_1$ is positive
semidefinite. Define the comparison and sensitivity constants
\begin{equation}
a_0=\min\{1,\underline w\},\quad
b_0=\max\{1,\overline w\},\quad
B=\overline wL^2,\quad C_F=L_AL^3,\quad
\gamma_0=\underline wL^2/d.
\label{eq:app-constants}
\end{equation}
For every $\theta\in\Theta$, the curvature bounds imply
\begin{align}
\rho I+\sum_{i=1}^nG_\theta(x_i)-a_0H_n
&=(1-a_0)\rho I+
\sum_{i=1}^n\big(w_\theta(x_i)-a_0\big)x_ix_i^\top\succeq0,\notag\\
b_0H_n-\rho I-\sum_{i=1}^nG_\theta(x_i)
&=(b_0-1)\rho I+
\sum_{i=1}^n\big(b_0-w_\theta(x_i)\big)x_ix_i^\top\succeq0.
\label{eq:app-matrix-comparison-proof}
\end{align}
In particular,
\begin{equation}
a_0H_n\preceq\rho I+\sum_{i=1}^nG_\theta(x_i)
\preceq b_0H_n,\qquad
0\preceq G_\theta(x)\preceq BI.
\label{eq:app-comparability}
\end{equation}
Convexity of $\Theta$ permits the same comparison for curvature
integrated along a segment between feasible parameters.
Since $\norm{xx^\top}_{\op}=\norm{x}_2^2$, information sensitivity is
\begin{align}
\norm{G_\theta(x)-G_{\theta'}(x)}_{\op}
&=|A''(x^\top\theta)-A''(x^\top\theta')|\norm{x}_2^2\notag\\
&\le L_A|x^\top(\theta-\theta')|\norm{x}_2^2
\le L_A\norm{x}_2^3\norm{\theta-\theta'}_2
\le C_F\norm{\theta-\theta'}_2.
\label{eq:app-fisher-lipschitz}
\end{align}
Constants in finite-time bounds depend only on
$d,L,S,\rho,\underline w,\overline w,L_A,\sigma,\tau_0$.
A dependence on $\gamma^\star$ can be replaced by its lower bound
$\gamma_0$ when an inverse power is used. Fourth-moment constants in
the fixed-$m$ efficiency proof may additionally depend on $m$.

\subsubsection*{Predictability without within-batch feedback.}
For the analysis only, reveal responses in flat-index order. Let
$\mathscr F_{i-1}$ contain the revealed responses preceding $i$ and
the entire already selected batch containing $i$. At a batch boundary,
include the new batch choice and any independent policy randomization
before revealing that batch's first response. These sigma-fields are
nested: a new batch is formed only from completed earlier data and
parameter-independent randomization. Thus $x_i$ is
$\mathscr F_{i-1}$-measurable. For $i=(t-1)m+j$, conditional
independence within $\mathcal B_t$ yields
\begin{align}
\E[\xi_i\mid\mathscr F_{i-1}]&=0,\notag\\
\E[\xi_i^2\mid\mathscr F_{i-1}]&=w_{\theta^\star}(x_i),\notag\\
\E[e^{a\xi_i}\mid\mathscr F_{i-1}]&\le e^{a^2\sigma^2/2}
\quad(a\in\R).
\label{eq:app-predictable-noise}
\end{align}
Earlier responses within the same batch do not change the next
response's conditional law given that batch and its preceding data.
This analytical revelation does not allow the policy to use those
responses when choosing its batch. In particular, we never condition
on the entire future, adaptively selected query sequence.

\subsubsection*{Score covariance.}
For $i<k$, $x_i\xi_i x_k^\top$ is
$\mathscr F_{k-1}$-measurable. Consequently,
\begin{equation}
\E[x_i\xi_i\xi_kx_k^\top]
=\E\!\left[x_i\xi_i x_k^\top
\E[\xi_k\mid\mathscr F_{k-1}]\right]=0.
\label{eq:app-cross-score-zero}
\end{equation}
The diagonal terms satisfy
\begin{equation}
\E[x_i\xi_i^2x_i^\top]
=\E[x_ix_i^\top\E[\xi_i^2\mid\mathscr F_{i-1}]]
=\E[G_{\theta^\star}(x_i)].
\label{eq:app-diagonal-score}
\end{equation}
Expanding the double sum in $U_nU_n^\top$ therefore gives
\begin{equation}
\E[U_n]=0,\qquad
\E[U_nU_n^\top]=\sum_{i=1}^n\E[G_{\theta^\star}(x_i)]
=\E[Z_n-\rho I].
\label{eq:app-score-cov}
\end{equation}
All expectations here are under $\theta^\star$.

\subsubsection*{Existence of the implemented optimizers.}
For fixed data, the objective in \eqref{eq:mle} is continuous on the
compact set $\Theta$, and its Hessian dominates $\rho I$. It therefore
has a unique minimizer. For a fixed temporary state and fitted
parameter, the query score is continuous on the compact ball; its
maximum is attained. The fixed measurable tie rule specified in the
algorithm selects an admissible maximizer. No approximate oracle is
substituted in the proof.

\subsection{Proofs of Spectral Smoothing Bounds}
\label{app:spectral}

For a matrix $K$, $\norm{K}_1$ denotes the sum of its singular values
(the trace norm); $\norm{K}_{\op}$ denotes its largest singular value.
For $1\le p<\infty$, the Schatten norm is
$\norm{K}_p=(\sum_r s_r(K)^p)^{1/p}$, where $s_r(K)$ are the singular
values. A density matrix is positive semidefinite with trace one.
If it is positive definite, its entropy is
$\mathcal H(P)=-\tr(P\log P)$.

\begin{lemma}[Spectral sensitivity]
\label{lem:gibbs}
For symmetric $A_1,A_2,D$ and $\tau>0$,
\begin{align}
\norm{P_\tau(A_1)-P_\tau(A_2)}_1
&\le 2\norm{A_1-A_2}_{\op}/\tau,
\label{eq:app-gibbs-lip}\\
\Phi_\tau(A_1+D)-\Phi_\tau(A_1)
&\ge\tr(P_\tau(A_1)D)-\norm{D}_{\op}^2/\tau.
\label{eq:app-potential-increment}
\end{align}
\end{lemma}

\begin{proof}
\textit{Matrix-exponential derivative.}
For square matrices $K,E$, let $\mathrm D\exp(K)[E]$ denote the
coefficient of the first-order variation of $\exp(K)$ in direction $E$.
Termwise differentiation of the absolutely convergent exponential
series on bounded matrix sets gives
\begin{equation}
\mathrm D\exp(K)[E]
=\sum_{r=1}^\infty\frac1{r!}
\sum_{a=0}^{r-1}K^aEK^{r-1-a}
=\int_0^1e^{sK}Ee^{(1-s)K}\,ds.
\label{eq:app-exp-derivative}
\end{equation}
The integral identity follows by multiplying the two series:
\begin{equation}
\int_0^1\frac{s^a(1-s)^b}{a!b!}\,ds
=\frac1{(a+b+1)!},\qquad a,b\in\{0,1,\ldots\}.
\label{eq:app-beta-integral}
\end{equation}
The integral evaluation follows by repeated integration by parts;
its value is $a!b!/(a+b+1)!$ before division by $a!b!$.
Thus no commutation between $K$ and $E$ is required.

Fix $A=A_1$. Define the matrix-valued map $F(A)=e^{-A/\tau}$
and its scalar normalizer $z(A)=\tr F(A)$, and abbreviate
$F=F(A)$, $z=z(A)$, and $P=F/z$. The symbols
$\mathrm DF(A)[D]$ and $\mathrm Dz(A)[D]$ denote their derivatives
in direction $D$. By \eqref{eq:app-exp-derivative} and cyclicity of the trace,
\begin{align}
\mathrm DF(A)[D]&=-\frac1\tau
\int_0^1e^{-sA/\tau}De^{-(1-s)A/\tau}\,ds,\notag\\
\mathrm Dz(A)[D]&=-\frac1\tau\tr(FD).
\label{eq:app-exp-normalizer}
\end{align}
Since $e^{-sA/\tau}=z^sP^s$, the quotient rule yields
\begin{align}
\mathrm DP_\tau(A)[D]
&=\frac{\mathrm DF(A)[D]}z-\frac{F\,\mathrm Dz(A)[D]}{z^2}\notag\\
&=-\frac1\tau\int_0^1P^sDP^{1-s}\,ds
+\frac1\tau P\tr(PD).
\label{eq:app-duhamel}
\end{align}
For $0<s<1$, the Schatten H\"older inequality gives
\begin{align}
\norm{P^sDP^{1-s}}_1
&\le\norm{P^s}_{1/s}\norm{D}_{\op}
\norm{P^{1-s}}_{1/(1-s)}\notag\\
&=(\tr P)^s\norm{D}_{\op}(\tr P)^{1-s}
=\norm{D}_{\op}.
\label{eq:app-holder-density}
\end{align}
The same bound at $s=0,1$ follows by continuity. Moreover,
\begin{equation}
\norm{P\tr(PD)}_1
=|\tr(PD)|\tr P\le\norm{D}_{\op}.
\label{eq:app-trace-part}
\end{equation}
It follows that
$\norm{\mathrm DP_\tau(A)[D]}_1\le2\norm{D}_{\op}/\tau$.
Applying this derivative bound on the segment joining $A_1$ and $A_2$
proves
\begin{align}
\norm{P_\tau(A_2)-P_\tau(A_1)}_1
&=\left\|\int_0^1\mathrm DP_\tau(A_1+s(A_2-A_1))
[A_2-A_1]\,ds\right\|_1\notag\\
&\le\int_0^1\frac{2\norm{A_2-A_1}_{\op}}\tau\,ds.
\label{eq:app-integrated-sensitivity}
\end{align}

\textit{Potential increment.}
Differentiating $\Phi_\tau(A)=-\tau\log z$ gives
\begin{equation}
\mathrm D\Phi_\tau(A)[D]
=-\tau\frac{\mathrm Dz(A)[D]}z=\tr(P_\tau(A)D).
\label{eq:app-potential-derivative}
\end{equation}
Using $|\tr(KD)|\le\norm{K}_1\norm{D}_{\op}$ and the first part,
\begin{align}
\Phi_\tau(A+D)-\Phi_\tau(A)-\tr(P_\tau(A)D)
&=\int_0^1\tr\!\big((P_\tau(A+sD)-P_\tau(A))D\big)\,ds\notag\\
&\ge-\int_0^1\frac{2s\norm{D}_{\op}^2}{\tau}\,ds
=-\frac{\norm{D}_{\op}^2}{\tau}.
\label{eq:app-smoothness-integral}
\end{align}
This proves the second inequality.
\end{proof}

\subsubsection*{Energy and temperature identities used below.}
Let $\lambda_r$ be the eigenvalues of a symmetric $A$, and let
$\lambda=\min_r\lambda_r$. Since
$1\le\sum_r e^{-(\lambda_r-\lambda)/\tau}\le d$,
\begin{equation}
\Phi_\tau(A)
=\lambda-\tau\log\!\left(\sum_{r=1}^d
 e^{-(\lambda_r-\lambda)/\tau}\right)
\in[\lambda-\tau\log d,\lambda].
\label{eq:app-soft-min-bracket}
\end{equation}
For the eigenvalues $p_r>0$ of $P=P_\tau(A)$,
$\sum_rp_r=1$ and $p_r\le1$, so $\mathcal H(P)\ge0$.
Concavity of the scalar logarithm gives
\begin{equation}
\mathcal H(P)=\sum_{r=1}^dp_r\log(1/p_r)
\le\log\!\left(\sum_{r=1}^dp_r/p_r\right)=\log d.
\label{eq:app-entropy-bound}
\end{equation}
The equality $\log P=-A/\tau-(\log z)I$ implies
\begin{equation}
\Phi_\tau(A)=\tr(PA)-\tau\mathcal H(P),\qquad
\tr(PA)\le\lambda_{\min}(A)+\tau\log d.
\label{eq:app-gibbs-energy}
\end{equation}
At fixed $A$,
$\partial_\tau z=\tau^{-2}\tr(Ae^{-A/\tau})$. Therefore
\begin{equation}
\partial_\tau\Phi_\tau(A)
=-\log z-\tau^{-1}\tr(PA)=-\mathcal H(P_\tau(A)).
\label{eq:app-temperature-derivative}
\end{equation}
For $\tau'\ge\tau$, integration gives
\begin{equation}
\Phi_{\tau'}(A)-\Phi_\tau(A)
=-\int_\tau^{\tau'}\mathcal H(P_v(A))\,dv
\ge-(\tau'-\tau)\log d.
\label{eq:app-temp-change}
\end{equation}

\subsection{Proofs of Deterministic Excitation}
\label{app:excitation}

\begin{lemma}[Input excitation]
\label{lem:excitation}
There are $c_E>0$ and $N_E\ge1$, independent of $m$, such that
\begin{equation}
\lambda_{\min}(H_N)\ge c_E\frac N{\log(e+N)}
\qquad(N=mT\ge N_E)
\label{eq:app-excitation}
\end{equation}
for the sample-scaled policy, even with arbitrary fitted parameters
in $\Theta$.
\end{lemma}

\begin{proof}
\textit{Step 1: every chosen query has nonzero weighted information.}
Consider query $i+1$ in round $t$, with $i$ previous queries selected,
and denote its temporary plug-in matrix by $W_i$.
All $i$ queries are evaluated at the current round's fitted parameter,
so \eqref{eq:app-comparability} gives
$a_0H_i\preceq W_i\preceq b_0H_i$.
Set $P=P_{\tau_t}(W_i)$ and let $v$ be a unit eigenvector for
$\lambda_{\max}(P)$. Since $\tr P=1$, $\lambda_{\max}(P)\ge1/d$.
The query $Lv$ belongs to the full ball. Exact score maximization gives
\begin{align}
\overline w\,x_{i+1}^\top Px_{i+1}
&\ge w_{\widehat\theta_{t-1}}(x_{i+1})x_{i+1}^\top Px_{i+1}\notag\\
&\ge w_{\widehat\theta_{t-1}}(Lv)L^2v^\top Pv
\ge\underline wL^2/d.
\label{eq:app-score-lower-proof}
\end{align}
Thus
\begin{equation}
x_{i+1}^\top Px_{i+1}
\ge\frac{\underline wL^2}{\overline w d}.
\label{eq:app-score-lower}
\end{equation}

\textit{Step 2: convert spectral score into input leverage.}
Define $q_i=x_{i+1}^\top H_i^{-1}x_{i+1}$ and the terminal minimum
eigenvalue $h_N=\lambda_{\min}(H_N)$.
The matrix $H_i^{1/2}PH_i^{1/2}$ is positive semidefinite. Its largest
eigenvalue is at most its trace, so
\begin{align}
x_{i+1}^\top Px_{i+1}
&=(H_i^{-1/2}x_{i+1})^\top
 H_i^{1/2}PH_i^{1/2}(H_i^{-1/2}x_{i+1})\notag\\
&\le\tr(PH_i)q_i.
\label{eq:app-score-leverage}
\end{align}
By matrix comparison and \eqref{eq:app-gibbs-energy},
\begin{align}
\tr(PH_i)
&\le a_0^{-1}\tr(PW_i)\notag\\
&\le a_0^{-1}\big(\lambda_{\min}(W_i)+\tau_t\log d\big)\notag\\
&\le a_0^{-1}\big(b_0\lambda_{\min}(H_i)+\tau_t\log d\big)\notag\\
&\le a_0^{-1}\big(b_0h_N+\tau_T\log d\big).
\label{eq:app-energy-bound}
\end{align}
The third line follows by evaluating $W_i\preceq b_0H_i$ at a minimum
eigenvector of $H_i$. The last line uses $H_i\preceq H_N$ and the
increasing temperature. Define
\begin{equation}
c_0=\frac{a_0\underline wL^2}{\overline w d},\qquad
q_{\max}=L^2/\rho,\qquad c_1=\frac{c_0}{1+q_{\max}}.
\label{eq:app-excitation-constants}
\end{equation}
Equations \eqref{eq:app-score-lower}--\eqref{eq:app-energy-bound}, and
$H_i\succeq\rho I$, yield
\begin{equation}
\frac{c_0}{b_0h_N+\tau_T\log d}\le q_i\le q_{\max}.
\label{eq:app-leverage-lower}
\end{equation}

\textit{Step 3: sum leverage through a determinant identity.}
Since $H_{i+1}=H_i+x_{i+1}x_{i+1}^\top$,
\begin{equation}
\det H_{i+1}=\det H_i\,
\det\!\left(I+H_i^{-1/2}x_{i+1}x_{i+1}^\top H_i^{-1/2}\right)
=\det H_i(1+q_i).
\label{eq:app-determinant-update}
\end{equation}
The last matrix has eigenvalues $1+q_i,1,\ldots,1$.
For $0\le q\le q_{\max}$,
\begin{equation}
\log(1+q)=\int_0^q\frac{dv}{1+v}
\ge\frac q{1+q_{\max}}.
\label{eq:app-log-leverage}
\end{equation}
If $\lambda_1(H_N),\ldots,\lambda_d(H_N)$ are its positive eigenvalues,
the arithmetic--geometric mean inequality gives
\begin{align}
\frac{\det H_N}{\det H_0}
&=\prod_{r=1}^d\frac{\lambda_r(H_N)}\rho
\le\left(\frac{\tr H_N}{\rho d}\right)^d
\le\left(1+\frac{NL^2}{\rho d}\right)^d.
\label{eq:app-determinant-upper}
\end{align}
Combining these displays proves the explicit inequality
\begin{equation}
\frac{Nc_1}{b_0h_N+\tau_T\log d}
\le\sum_{i=0}^{N-1}\log(1+q_i)
=\log\frac{\det H_N}{\det H_0}
\le d\log\!\left(1+\frac{NL^2}{\rho d}\right).
\label{eq:app-excitation-det}
\end{equation}
This determinant is only an analytical device; the policy has not
been changed to a log-determinant objective.

\textit{Step 4: give constants and an $m$-independent threshold.}
Define
\begin{equation}
a=L^2/(\rho d),\qquad c_h=1+\log\max\{1,a\}.
\label{eq:app-log-constant}
\end{equation}
For $N\ge1$, $1+aN\le\max\{1,a\}(1+N)$ and $\log(e+N)\ge1$ imply
$\log(1+aN)\le c_h\log(e+N)$. Rearranging
\eqref{eq:app-excitation-det} with $\tau_T=\tau_0N^{3/4}$ gives
\begin{align}
h_N
&\ge\frac1{b_0}\left[
\frac{Nc_1}{d\log(1+aN)}-\tau_0N^{3/4}\log d\right]\notag\\
&\ge\frac{c_1N}{b_0dc_h\log(e+N)}
-\frac{\tau_0\log d}{b_0}N^{3/4}.
\label{eq:app-explicit-excitation}
\end{align}
Because $\log(e+N)/N^{1/4}\to0$, choose $N_E$ so that, for all
integers $N\ge N_E$,
\begin{equation}
\frac{\log(e+N)}{N^{1/4}}
\le\frac{c_1}{2dc_h\tau_0\log d}.
\label{eq:app-excitation-threshold}
\end{equation}
This choice depends only on the displayed model constants, not on $m$.
The negative term in \eqref{eq:app-explicit-excitation} is then at most
half the positive term. Set
\begin{equation}
c_E=\frac{c_1}{2b_0dc_h}.
\label{eq:app-excitation-final-constant}
\end{equation}
This proves the lemma. By \eqref{eq:app-comparability}, both
$V_T^\star$ and $\Vhat_T$ have minimum eigenvalue at least $a_0h_N$.
No concentration inequality, parameter consistency, or accurate
initial fit was used.
\end{proof}

\subsection{Proofs of Confidence and Preliminary Localization}
\label{app:localization}

For $n\ge1$ and $\delta\in(0,1)$, define the score radius
\begin{equation}
b_n(\delta)=\sigma\sqrt{
 d\log\!\left(1+\frac{nL^2}{\rho d}\right)+2\log\frac2\delta}.
\label{eq:app-score-radius}
\end{equation}
We first derive its uniform-in-time probability guarantee, then use
constrained optimality. This is the self-normalized method of
\citet[Theorem 1]{AbbasiYadkori2011}, written here for the precise
filtration and ridge matrix used in this manuscript.

\begin{lemma}[Time-uniform score event]
\label{lem:score}
There is an event $\mathcal E_\delta$ with
$\Prob(\mathcal E_\delta)\ge1-\delta/2$ on which
\begin{equation}
\norm{U_n}_{H_n^{-1}}\le b_n(\delta)\qquad\text{for all }n\ge1.
\label{eq:app-self-normalized}
\end{equation}
\end{lemma}

\begin{proof}
\textit{Step 1: exponential supermartingales.}
The positive conditional response variance ensures $\sigma>0$.
For a fixed vector $z\in\R^d$, define
\begin{equation}
\mathcal M_n(z)=\exp\!\left\{
 \frac{z^\top U_n}{\sigma}
 -\frac12\sum_{i=1}^n(z^\top x_i)^2\right\}.
\label{eq:app-fixed-exponential}
\end{equation}
Using \eqref{eq:app-predictable-noise},
\begin{align}
\E[\mathcal M_n(z)\mid\mathscr F_{n-1}]
&=\mathcal M_{n-1}(z)
\E\!\left[\exp\!\left\{\frac{z^\top x_n}{\sigma}\xi_n
-\frac12(z^\top x_n)^2\right\}\middle|\mathscr F_{n-1}\right]
\le\mathcal M_{n-1}(z).
\label{eq:app-supermartingale-step}
\end{align}
Thus this is a nonnegative supermartingale with initial value one.

\textit{Step 2: compute a Gaussian mixture exactly.}
Let $q_\rho(z)=(\rho/(2\pi))^{d/2}e^{-\rho\norm{z}_2^2/2}$ be the
Gaussian density with covariance $\rho^{-1}I$, and define
$\overline{\mathcal M}_n=\int\mathcal M_n(z)q_\rho(z)\,dz$.
Nonnegativity permits conditional integration of
\eqref{eq:app-supermartingale-step}, so the mixture is also a
nonnegative supermartingale with $\overline{\mathcal M}_0=1$.
The exponent in the integral satisfies
\begin{align}
\frac{z^\top U_n}{\sigma}-\frac12z^\top H_nz
&=-\frac12\left(z-\frac{H_n^{-1}U_n}{\sigma}\right)^\top
 H_n\left(z-\frac{H_n^{-1}U_n}{\sigma}\right)
 +\frac{\norm{U_n}_{H_n^{-1}}^2}{2\sigma^2}.
\label{eq:app-mixture-square}
\end{align}
Translation and diagonalization of $H_n$ give
\begin{align}
\overline{\mathcal M}_n
&=\left(\frac\rho{2\pi}\right)^{d/2}
 e^{\norm{U_n}_{H_n^{-1}}^2/(2\sigma^2)}
 \int_{\R^d}e^{-z^\top H_nz/2}\,dz\notag\\
&=\frac{\det(\rho I)^{1/2}}{\det(H_n)^{1/2}}
 e^{\norm{U_n}_{H_n^{-1}}^2/(2\sigma^2)}.
\label{eq:app-mixture-evaluation}
\end{align}
Indeed, in an eigenbasis, the integral is the product of
$\int_\R e^{-\lambda_r(H_n)v^2/2}\,dv
=(2\pi/\lambda_r(H_n))^{1/2}$.

\textit{Step 3: control all times with one event.}
Let $\nu=\inf\{n\ge0:\overline{\mathcal M}_n\ge2/\delta\}$, with
$\inf\varnothing=\infty$. For an integer $k\ge1$,
\begin{equation}
\overline{\mathcal M}_{\nu\wedge k}
=1+\sum_{n=1}^k\mathbf1_{\{\nu\ge n\}}
(\overline{\mathcal M}_n-\overline{\mathcal M}_{n-1}).
\label{eq:app-stopped-sum}
\end{equation}
The indicator is $\mathscr F_{n-1}$-measurable. Taking conditional
expectations term by term shows
$\E[\overline{\mathcal M}_{\nu\wedge k}]\le1$.
Nonnegativity and the definition of $\nu$ imply
\begin{equation}
\frac2\delta\Prob(\nu\le k)
\le\E[\overline{\mathcal M}_{\nu\wedge k}]\le1.
\label{eq:app-stopped-threshold}
\end{equation}
Letting $k\to\infty$ gives $\Prob(\nu<\infty)\le\delta/2$.
Define $\mathcal E_\delta=\{\nu=\infty\}$. On this event,
\eqref{eq:app-mixture-evaluation} implies, for every $n$,
\begin{equation}
\norm{U_n}_{H_n^{-1}}^2
\le\sigma^2\left[
\log\frac{\det H_n}{\det(\rho I)}+2\log\frac2\delta\right].
\label{eq:app-determinant-confidence}
\end{equation}
The determinant estimate \eqref{eq:app-determinant-upper}, with $n$
in place of $N$, yields \eqref{eq:app-self-normalized}.
\end{proof}

For the ridge likelihood at a completed count, define
\begin{equation}
\mathcal L_N(\theta)=\sum_{i=1}^N
\big[A(x_i^\top\theta)-Y_ix_i^\top\theta\big]
+\frac\rho2\norm{\theta}_2^2.
\label{eq:app-loss-definition}
\end{equation}
For $e_T=\thetahat_T-\theta^\star$, its integrated Hessian is
\begin{equation}
K_N=\int_0^1\nabla^2\mathcal L_N(\theta^\star+s e_T)\,ds
=\rho I+\sum_{i=1}^N
\left[\int_0^1w_{\theta^\star+s e_T}(x_i)\,ds\right]x_ix_i^\top.
\label{eq:app-integrated-hessian}
\end{equation}
The segment is feasible by convexity, and hence
$a_0H_N\preceq K_N\preceq b_0H_N$.

\begin{lemma}[Constrained localization]
\label{lem:localization}
On $\mathcal E_\delta$, for every $N=mT$,
\begin{equation}
\norm{e_T}_{H_N}\le\frac{b_N(\delta)+\sqrt\rho S}{a_0},
\qquad\norm{e_T}_{\Vhat_T}\le\beta_N(\delta).
\label{eq:app-confidence}
\end{equation}
Moreover, $\norm{e_T}_2\le C_L\log(eN/\delta)/\sqrt N$ for a constant
$C_L$ independent of $m,T,\delta$.
\end{lemma}

\begin{proof}
\textit{Step 1: use the constrained inequality, not a zero gradient.}
Differentiation gives
\begin{align}
\nabla\mathcal L_N(\theta^\star)&=-U_N+\rho\theta^\star,\notag\\
\nabla\mathcal L_N(\thetahat_T)-\nabla\mathcal L_N(\theta^\star)
&=\int_0^1\nabla^2\mathcal L_N(\theta^\star+s e_T)e_T\,ds
=K_Ne_T.
\label{eq:app-gradient-expansion}
\end{align}
Since $\thetahat_T$ minimizes over the convex set $\Theta$,
$\nabla\mathcal L_N(\thetahat_T)^\top(\theta^\star-\thetahat_T)\ge0$.
Substituting \eqref{eq:app-gradient-expansion} gives
\begin{equation}
e_T^\top K_Ne_T
\le e_T^\top(U_N-\rho\theta^\star).
\label{eq:app-variational}
\end{equation}
Thus a boundary solution causes no difficulty.

\textit{Step 2: obtain the ellipsoid.}
Cauchy--Schwarz in the $H_N$ norm gives
\begin{align}
a_0\norm{e_T}_{H_N}^2
&\le e_T^\top K_Ne_T
\le\norm{e_T}_{H_N}\norm{U_N-\rho\theta^\star}_{H_N^{-1}}\notag\\
&\le\norm{e_T}_{H_N}
\left(\norm{U_N}_{H_N^{-1}}+
\rho\norm{\theta^\star}_{H_N^{-1}}\right)\notag\\
&\le\norm{e_T}_{H_N}\big(b_N(\delta)+\sqrt\rho S\big).
\label{eq:app-constrained-cauchy}
\end{align}
The last step uses $H_N^{-1}\preceq\rho^{-1}I$ and
\cref{lem:score}. Division by $\norm{e_T}_{H_N}$ is valid when $e_T\ne0$;
when $e_T=0$, the bound is immediate. Since
$\Vhat_T\preceq b_0H_N$,
\begin{equation}
\norm{e_T}_{\Vhat_T}
\le\sqrt{b_0}\norm{e_T}_{H_N}
\le\frac{\sqrt{b_0}}{a_0}\big(b_N(\delta)+\sqrt\rho S\big)
=\beta_N(\delta).
\label{eq:app-ellipsoid-result}
\end{equation}

\textit{Step 3: make preliminary localization uniform in $m$.}
Define $\Lambda_N=\log(eN/\delta)\ge1$ and, using $c_h$ from
\eqref{eq:app-log-constant}, set
\begin{equation}
K_b=\sigma\sqrt{dc_h+2+2\log2}+\sqrt\rho S.
\label{eq:app-radius-constant}
\end{equation}
The inequalities $\log(1+aN)\le c_h\Lambda_N$ and
$2\log(2/\delta)\le(2+2\log2)\Lambda_N$ imply
\begin{equation}
b_N(\delta)+\sqrt\rho S\le K_b\sqrt{\Lambda_N}.
\label{eq:app-radius-log-control}
\end{equation}
For $N\ge N_E$, combine \cref{lem:excitation} with
$\log(e+N)\le2\Lambda_N$ to obtain
\begin{align}
\norm{e_T}_2
&\le\frac{\norm{e_T}_{H_N}}{\sqrt{\lambda_{\min}(H_N)}}
\le\frac{K_b}{a_0}\sqrt{\Lambda_N}
\sqrt{\frac{\log(e+N)}{c_EN}}\notag\\
&\le\frac{\sqrt2K_b}{a_0\sqrt{c_E}}
\frac{\Lambda_N}{\sqrt N}.
\label{eq:app-large-count-localization}
\end{align}
For $1\le N<N_E$, compactness gives
$\norm{e_T}_2\le2S\le2S\sqrt{N_E}\Lambda_N/\sqrt N$.
Therefore the explicit choice
\begin{equation}
C_L=\max\left\{\frac{\sqrt2K_b}{a_0\sqrt{c_E}},\,
2S\sqrt{N_E}\right\}
\label{eq:app-localization-constant}
\end{equation}
works for every completed count, proving
\begin{equation}
\norm{e_T}_2\le C_L\frac{\log(eN/\delta)}{\sqrt N}.
\label{eq:app-preliminary-localization}
\end{equation}
Neither \cref{thm:information} nor optimal information growth was used.
\end{proof}

\subsection{Proofs of Proposition 1 and the Design Benchmark}
\label{app:ceiling}

\begin{proof}[Proof of \cref{prop:ceiling}]
\textit{Step 1: compactness and attainment.}
Let $\mathbb S^d$ be the real symmetric $d\times d$ matrices, a vector
space of dimension $q=d(d+1)/2$. Define
\begin{equation}
\mathcal K=\{G_{\theta^\star}(x):x\in\mathcal X\},\qquad
\mathcal C=\{M(\nu):\nu\in\mathcal P(\mathcal X)\}.
\label{eq:app-feasible-matrices}
\end{equation}
Continuity of $G_{\theta^\star}$ and compactness of the ball make
$\mathcal K$ compact. We spell out why its convex hull is compact.
Any convex combination of more than $q+1$ matrices can be shortened:
if its positive weights are $p_r$ and matrices are $A_r\in\mathcal K$,
the vectors $(A_r,1)$ are linearly dependent. Thus there are coefficients
$\alpha_r$, not all zero, with
$\sum_r\alpha_rA_r=0$ and $\sum_r\alpha_r=0$.
There are both positive and negative $\alpha_r$. Taking
$s=\min_{\alpha_r>0}p_r/\alpha_r$ replaces the weights by
$p_r-s\alpha_r\ge0$, preserves their sum and weighted matrix, and
sets at least one weight to zero. Repeating reduces to $q+1$ terms.
Consequently the convex hull is the continuous image of the compact
product of $q+1$ copies of $\mathcal K$ and their probability simplex,
and is compact.

Every finite convex combination is $M(\nu)$ for a finite-support
measure. Conversely, uniform continuity of $G_{\theta^\star}$ on
$\mathcal X$ allows its integral under any probability measure to be
approximated, in operator norm, by finite weighted sums: partition the
ball into finitely many sets on which $G$ varies by at most $\epsilon$,
choose one point from each nonempty set, and use the measure of that
set as its weight. The approximation error is at most $\epsilon$.
The convex hull is closed, so the integral also belongs to it.
Thus $\mathcal C$ equals this compact convex hull.

For symmetric $A_1,A_2$, the Rayleigh representation gives
\begin{equation}
|\lambda_{\min}(A_1)-\lambda_{\min}(A_2)|
\le\norm{A_1-A_2}_{\op}.
\label{eq:app-rayleigh-continuity}
\end{equation}
For example, evaluating $A_1$ at a minimum unit eigenvector of $A_2$
gives one side, and exchanging the two matrices gives the other.
Hence $\lambda_{\min}$ is continuous on $\mathcal C$ and attains its
maximum, $\gamma^\star$.

\textit{Step 2: bound the optimal rate.}
Let $e_1,\ldots,e_d$ be the standard coordinate unit vectors and let
$\nu_0$ give weight $1/d$ to each query $Le_r$. Then
\begin{equation}
M(\nu_0)=\frac{L^2}{d}\sum_{r=1}^d
 w_{\theta^\star}(Le_r)e_re_r^\top
\succeq\frac{\underline wL^2}{d}I.
\label{eq:app-coordinate-design}
\end{equation}
For every feasible $\nu$,
\begin{align}
d\lambda_{\min}(M(\nu))
&\le\tr(M(\nu))
=\int w_{\theta^\star}(x)\norm{x}_2^2\,d\nu(x)
\le\overline wL^2=B.
\label{eq:app-trace-rate-bound}
\end{align}
Taking the maximum gives \eqref{eq:rate-range}.

\textit{Step 3: apply the feasible-set bound pathwise.}
For any realized adaptive query sequence, define its empirical
probability measure by $\nu_N=N^{-1}\sum_{i=1}^N\delta_{x_i}$,
where $\delta_x$ denotes unit mass at $x$. Then
\begin{equation}
\frac{Z_N-\rho I}{N}=M(\nu_N),\qquad
\lambda_{\min}(Z_N)=\rho+N\lambda_{\min}(M(\nu_N))
\le\rho+N\gamma^\star.
\label{eq:app-pathwise-ceiling}
\end{equation}
This holds on every path and does not require expectations.
\end{proof}

\subsubsection*{The ideal score needed for the policy comparison.}
For any density matrix $P$ and any maximizing measure $\nu^\star$,
$M(\nu^\star)\succeq\gamma^\star I$. Therefore
\begin{align}
\max_{x\in\mathcal X}\tr(PG_{\theta^\star}(x))
&\ge\int\tr(PG_{\theta^\star}(x))\,d\nu^\star(x)\notag\\
&=\tr(PM(\nu^\star))
\ge\gamma^\star\tr P=\gamma^\star.
\label{eq:app-ideal-score}
\end{align}
The last trace inequality follows because
$P^{1/2}(M(\nu^\star)-\gamma^\star I)P^{1/2}$ is positive semidefinite.
No interchange of a minimum and maximum is needed.

\subsection{Proofs of Theorem 1 and the Batched Gain Bound}
\label{app:information}

\subsubsection{Temporary states and score comparison}

Let $n_{t,j}=(t-1)m+j-1$ count the queries already selected before
position $j$ of batch $t$. Define the true and plug-in temporary states by
\begin{align}
Z_{t,j-1}
&=V_{t-1}^\star+\sum_{k=1}^{j-1}G_{\theta^\star}(x_{t,k}),\notag\\
W_{t,j-1}
&=\rho I+\sum_{s=1}^{t-1}\sum_{k=1}^m
 G_{\widehat\theta_{t-1}}(x_{s,k})
 +\sum_{k=1}^{j-1}G_{\widehat\theta_{t-1}}(x_{t,k}).
\label{eq:app-true-temporary}
\end{align}
Thus $Z_{t,j-1}=Z_{n_{t,j}}$, while $W_{t,j-1}$ is exactly the
state used by the algorithm. Both evaluate the same query prefix.
Write $\eta_t=\norm{\widehat\theta_t-\theta^\star}_2$, including
$\eta_0\le2S$. Define the true and plug-in scores for a candidate $x$ by
\begin{equation}
s_{t,j}^\star(x)=\tr(P_{\tau_t}(Z_{t,j-1})G_{\theta^\star}(x)),
\qquad
\widehat s_{t,j}(x)=\tr(P_{\tau_t}(W_{t,j-1})
 G_{\widehat\theta_{t-1}}(x)).
\label{eq:app-two-scores}
\end{equation}
Define a score-discrepancy bound and a bound on optimization loss by
\begin{equation}
d_{t,j}=C_F\eta_{t-1}\left(1+\frac{2Bn_{t,j}}{\tau_t}\right),
\qquad \ell_{t,j}=\min\{B,2d_{t,j}\}.
\label{eq:app-score-loss}
\end{equation}
These quantities are for analysis only.

\begin{lemma}[Greedy batch gain]
\label{lem:batch}
For every batch and every realized history,
\begin{equation}
\Phi_{\tau_t}(V_t^\star)-\Phi_{\tau_t}(V_{t-1}^\star)
\ge m\gamma^\star-\frac{mB^2}{\tau_t}
-\sum_{j=1}^m\ell_{t,j}.
\label{eq:app-batch-gain}
\end{equation}
\end{lemma}

\begin{proof}
\textit{Step 1: control re-evaluation of the entire prefix.}
Apply \eqref{eq:app-fisher-lipschitz} to each of the $n_{t,j}$ matrices
in \eqref{eq:app-true-temporary}:
\begin{align}
\norm{W_{t,j-1}-Z_{t,j-1}}_{\op}
&\le\sum_{i=1}^{n_{t,j}}
\norm{G_{\widehat\theta_{t-1}}(x_i)-G_{\theta^\star}(x_i)}_{\op}
\le n_{t,j}C_F\eta_{t-1}.
\label{eq:app-state-error}
\end{align}
For brevity in the next display write
$\widehat P=P_{\tau_t}(W_{t,j-1})$ and
$P^\star=P_{\tau_t}(Z_{t,j-1})$.
Adding and subtracting $\tr(\widehat P G_{\theta^\star}(x))$ gives
\begin{align}
|\widehat s_{t,j}(x)-s_{t,j}^\star(x)|
&=\left|\tr\!\left(\widehat P
 [G_{\widehat\theta_{t-1}}(x)-G_{\theta^\star}(x)]\right)
 +\tr\!\left((\widehat P-P^\star)G_{\theta^\star}(x)\right)\right|\notag\\
&\le\norm{\widehat P}_1
\norm{G_{\widehat\theta_{t-1}}(x)-G_{\theta^\star}(x)}_{\op}
 +\norm{\widehat P-P^\star}_1\norm{G_{\theta^\star}(x)}_{\op}\notag\\
&\le C_F\eta_{t-1}
 +\frac{2B}{\tau_t}\norm{W_{t,j-1}-Z_{t,j-1}}_{\op}
\le d_{t,j}.
\label{eq:app-score-comparison}
\end{align}
We used $\norm{\widehat P}_1=\tr\widehat P=1$ and
\cref{lem:gibbs}. The factor $n_{t,j}/\tau_t$ is essential: the
weight error comes from the whole recomputed state, not just the
newest query.

\textit{Step 2: transfer exact maximization to the true score.}
Let $x_{t,j}^\star$ maximize $s_{t,j}^\star$ over the ball.
The implemented query maximizes $\widehat s_{t,j}$, hence
\begin{align}
s_{t,j}^\star(x_{t,j}^\star)-s_{t,j}^\star(x_{t,j})
&=[s_{t,j}^\star(x_{t,j}^\star)-\widehat s_{t,j}(x_{t,j}^\star)]
 +[\widehat s_{t,j}(x_{t,j}^\star)-\widehat s_{t,j}(x_{t,j})]\notag\\
&\hspace{1em}+[\widehat s_{t,j}(x_{t,j})-s_{t,j}^\star(x_{t,j})]
\le d_{t,j}+0+d_{t,j}.
\label{eq:app-argmax-transfer}
\end{align}
Every true score lies in $[0,B]$, so the same nonnegative difference
is at most $B$. Equation \eqref{eq:app-ideal-score} now yields
\begin{equation}
s_{t,j}^\star(x_{t,j})
\ge s_{t,j}^\star(x_{t,j}^\star)-\ell_{t,j}
\ge\gamma^\star-\ell_{t,j}.
\label{eq:app-selected-true-score}
\end{equation}

\textit{Step 3: sum single-query increments within the batch.}
Write $Z_{t,j}=Z_{t,j-1}+G_{\theta^\star}(x_{t,j})$.
By \eqref{eq:app-potential-increment},
\begin{align}
\Phi_{\tau_t}(Z_{t,j})-\Phi_{\tau_t}(Z_{t,j-1})
&\ge s_{t,j}^\star(x_{t,j})
 -\frac{\norm{G_{\theta^\star}(x_{t,j})}_{\op}^2}{\tau_t}\notag\\
&\ge\gamma^\star-\ell_{t,j}-B^2/\tau_t.
\label{eq:app-query-potential-gain}
\end{align}
Summing over $j=1,\ldots,m$, and using
$Z_{t,0}=V_{t-1}^\star$ and $Z_{t,m}=V_t^\star$, proves the lemma.
\end{proof}

The virtual update matters exactly in
\eqref{eq:app-query-potential-gain}: the comparison uses the refreshed
state before each query. One single application of smoothness to the
aggregate increment would instead use
\begin{equation}
\frac1{\tau_t}\left\|\sum_{j=1}^mG_{\theta^\star}(x_{t,j})\right\|_{\op}^2
\le\frac{(mB)^2}{\tau_t},
\label{eq:app-coarse-batch-remainder}
\end{equation}
rather than summing $m$ terms of size $B^2/\tau_t$.
This comparison of analytical remainders is not a performance lower
bound on other batch algorithms.

\subsubsection{Telescoping across batches and temperatures}

\begin{proof}[Proof of \cref{thm:information}]
Work on $\mathcal E_\delta$ from \cref{lem:score}. Its probability is
at least $1-\delta/2$, and it is the same event used throughout.
Fix a terminal count $N=mT$, and write $\Lambda_N=\log(eN/\delta)$.

\textit{Step 1: keep the temperature-change terms explicit.}
The exact telescoping identity is
\begin{align}
\Phi_{\tau_T}(V_T^\star)
={}&\Phi_{\tau_1}(V_0^\star)
 +\sum_{t=1}^T
 [\Phi_{\tau_t}(V_t^\star)-\Phi_{\tau_t}(V_{t-1}^\star)]\notag\\
&+\sum_{t=1}^{T-1}
 [\Phi_{\tau_{t+1}}(V_t^\star)-\Phi_{\tau_t}(V_t^\star)].
\label{eq:app-temperature-telescoping}
\end{align}
The first term is $\rho-\tau_1\log d$ because $V_0^\star=\rho I$.
For each term in the second sum, \eqref{eq:app-temp-change} gives
\begin{equation}
\Phi_{\tau_{t+1}}(V_t^\star)-\Phi_{\tau_t}(V_t^\star)
\ge-(\tau_{t+1}-\tau_t)\log d.
\label{eq:app-temperature-loss}
\end{equation}
Thus \cref{lem:batch} and
$\lambda_{\min}(Z_N)\ge\Phi_{\tau_T}(Z_N)$ imply
\begin{align}
\lambda_{\min}(Z_N)
&\ge\rho-\tau_1\log d+N\gamma^\star
 -B^2m\sum_{t=1}^T\frac1{\tau_t}
 -\sum_{t=1}^T\sum_{j=1}^m\ell_{t,j}
 -(\tau_T-\tau_1)\log d\notag\\
&=\rho+N\gamma^\star-\tau_T\log d
 -B^2m\sum_{t=1}^T\frac1{\tau_t}
 -\sum_{t=1}^T\sum_{j=1}^m\ell_{t,j}.
\label{eq:app-master-gap}
\end{align}
Only the true states $Z_n$ are additive. No telescoping identity is
asserted for the recomputed matrices $\Vhat_t$.

\textit{Step 2: evaluate every required power sum.}
For $0<p<1$ and $T\ge1$, monotonicity of $s^{-p}$ gives
\begin{align}
\sum_{t=1}^Tt^{-p}
&\le1+\int_1^T s^{-p}\,ds
=\frac{T^{1-p}-p}{1-p}
\le\frac{T^{1-p}}{1-p},\notag\\
m\sum_{t=1}^T(mt)^{-p}
&\le\frac{m^{1-p}T^{1-p}}{1-p}
=\frac{N^{1-p}}{1-p}.
\label{eq:app-power-sums}
\end{align}
In particular,
\begin{equation}
\tau_T=\tau_0N^{3/4},\qquad
m\sum_{t=1}^T\frac1{\tau_t}
=\frac1{\tau_0}m\sum_{t=1}^T(mt)^{-3/4}
\le\frac4{\tau_0}N^{1/4}.
\label{eq:app-temperature-sums}
\end{equation}

\textit{Step 3: bound delayed parameter errors and their sum.}
The first batch has $\sum_j\ell_{1,j}\le Bm$ without using any
localization. For $t\ge2$, there are $(t-1)m$ observed labels, and
\cref{lem:localization} gives
\begin{align}
\eta_{t-1}
&\le C_L\frac{\log(e(t-1)m/\delta)}{\sqrt{(t-1)m}}
\le\sqrt2C_L\frac{\Lambda_N}{\sqrt{mt}},\notag\\
\frac{n_{t,j}}{\tau_t}
&\le\frac{mt}{\tau_0(mt)^{3/4}}
=\tau_0^{-1}(mt)^{1/4}.
\label{eq:app-delayed-error-components}
\end{align}
The first inequality uses only completed prior batches; the comparison
$t-1\ge t/2$ explains the factor $\sqrt2$.
Define $C_s=2\sqrt2C_FC_L$. Substitution into
\eqref{eq:app-score-loss} yields
\begin{equation}
\ell_{t,j}\le C_s\Lambda_N
\left[(mt)^{-1/2}+\frac{2B}{\tau_0}(mt)^{-1/4}\right]
\qquad(t\ge2).
\label{eq:app-delayed-error}
\end{equation}
Sum this bound using \eqref{eq:app-power-sums} with $p=1/2,1/4$:
\begin{align}
\sum_{t=1}^T\sum_{j=1}^m\ell_{t,j}
&\le Bm+C_s\Lambda_N\left[
 m\sum_{t=2}^T(mt)^{-1/2}
 +\frac{2B}{\tau_0}m\sum_{t=2}^T(mt)^{-1/4}\right]\notag\\
&\le Bm+2C_s\Lambda_NN^{1/2}
 +\frac{8BC_s}{3\tau_0}\Lambda_NN^{3/4}.
\label{eq:app-total-plugin-loss}
\end{align}
For $T=1$ the sums starting at $2$ are empty and the same upper
bound remains valid.

\textit{Step 4: obtain a nonasymptotic true-information gap.}
The ceiling in \cref{prop:ceiling} and the previous calculations give
\begin{align}
0\le\rho+N\gamma^\star-\lambda_{\min}(Z_N)
\le{}&\tau_0\log d\,N^{3/4}
 +\frac{4B^2}{\tau_0}N^{1/4}+Bm\notag\\
&+2C_s\Lambda_NN^{1/2}
 +\frac{8BC_s}{3\tau_0}\Lambda_NN^{3/4}.
\label{eq:app-decomposed-gap}
\end{align}
Every coefficient in this display is independent of $m$.
The $Bm$ term is an upper bound for learning before the first labels,
not a lower bound on batch delay.

\textit{Step 5: compare the terminal recomputed information.}
By \eqref{eq:app-fisher-lipschitz} and preliminary localization,
\begin{align}
\norm{\Vhat_T-Z_N}_{\op}
&\le\sum_{i=1}^N
\norm{G_{\widehat\theta_T}(x_i)-G_{\theta^\star}(x_i)}_{\op}
\le NC_F\eta_T
\le C_FC_L\Lambda_N\sqrt N.
\label{eq:app-initial-tracking}
\end{align}
Define $D_N^{\mathrm{info}}=\rho+N\gamma^\star-\lambda_{\min}(Z_N)$,
a nonnegative scalar. The eigenvalue continuity estimate
\eqref{eq:app-rayleigh-continuity} gives
\begin{align}
\left|\frac{\lambda_{\min}(\Vhat_T)}N-\gamma^\star\right|
&\le\frac{\norm{\Vhat_T-Z_N}_{\op}}N
 +\left|\frac{\rho-D_N^{\mathrm{info}}}{N}\right|\notag\\
&\le\frac{\norm{\Vhat_T-Z_N}_{\op}}N
 +\frac\rho N+\frac{D_N^{\mathrm{info}}}N.
\label{eq:app-plugin-eigenvalue-bound}
\end{align}
For $N\ge1$ and $\Lambda_N\ge1$, each of
$N^{-1},N^{-3/4},N^{-1/2}$ is at most $\Lambda_NN^{-1/4}$.
For example, the explicit constant
\begin{equation}
C=1+\rho+B+\tau_0\log d+\frac{4B^2}{\tau_0}
 +2C_s+\frac{8BC_s}{3\tau_0}+C_FC_L
\label{eq:app-main-constant}
\end{equation}
therefore makes both inequalities in \eqref{eq:information-rate}
follow from \eqref{eq:app-decomposed-gap} and
\eqref{eq:app-plugin-eigenvalue-bound}.
All steps were deterministic on the single event $\mathcal E_\delta$,
so they hold for all completed horizons simultaneously. Since
$\Prob(\mathcal E_\delta)\ge1-\delta/2\ge1-\delta$, this proves the
formal theorem.
\end{proof}

\subsubsection{Tracking, probability-one limits, and known curvature}

\subsubsection*{Refined tracking on the same event.}
Suppose $\varepsilon_{N,m}(\delta)\le\gamma^\star/2$.
The second inequality in \eqref{eq:information-rate} implies
\begin{equation}
\lambda_{\min}(Z_N)
\ge\rho+N(\gamma^\star-\varepsilon_{N,m}(\delta))
\ge N\gamma^\star/2.
\label{eq:app-refined-information}
\end{equation}
Because $Z_N\preceq b_0H_N$,
$\lambda_{\min}(H_N)\ge N\gamma^\star/(2b_0)$.
Using \eqref{eq:app-confidence} and
\eqref{eq:app-radius-log-control},
\begin{equation}
\eta_T\le\frac{K_b}{a_0}
\sqrt{\frac{2b_0\Lambda_N}{N\gamma^\star}}
\le C_R\sqrt{\frac{\Lambda_N}{N}},\qquad
C_R=\frac{K_b}{a_0}\sqrt{\frac{2b_0}{\gamma_0}}.
\label{eq:app-refined-localization}
\end{equation}
The inequality $\norm{\Vhat_T-Z_N}_{\op}\le NC_F\eta_T$ then proves
\eqref{eq:tracking}, for example with $C'=1+C_FC_R$.

\subsubsection*{Almost-sure conclusions without a missing event argument.}
Fix $m$. For every fixed $\delta$, the bounds on $\mathcal E_\delta$
imply, as $T\to\infty$,
\begin{equation}
\varepsilon_{m,T}(\delta)\to0,\qquad
\frac{\log(emT/\delta)}{\sqrt{mT}}\to0.
\label{eq:app-deterministic-rate-limits}
\end{equation}
Let $\mathcal A_m$ be the event that both normalized minimum eigenvalues
converge to $\gamma^\star$ and that $\widehat\theta_T\to\theta^\star$.
The bounds imply $\mathcal E_\delta\subseteq\mathcal A_m$.
Consequently, for every integer $k\ge1$,
\begin{equation}
\Prob(\mathcal A_m^c)
\le\Prob(\mathcal E_{2^{-k}}^c)\le2^{-k-1},
\qquad\text{and hence}\qquad\Prob(\mathcal A_m)=1.
\label{eq:app-probability-one}
\end{equation}
Multiplication by the fixed $m$ changes the per-query limit
$\gamma^\star$ into the per-round limit $\gstar$ in
\eqref{eq:main-limit}. This argument does not assert simultaneous
limits for a batch size growing with $T$.

\subsubsection*{Parameter-independent curvature.}
If $C_F=0$, \eqref{eq:app-fisher-lipschitz} shows that the information
is independent of the fitted parameter. Then $\ell_{t,j}=0$ for all
queries and $\Vhat_T=Z_N$. Substituting into
\eqref{eq:app-master-gap} gives the deterministic refinement
\begin{equation}
0\le\frac{\rho+N\gamma^\star-\lambda_{\min}(V_T^\star)}N
\le\tau_0\log d\,N^{-1/4}
+\frac{4B^2}{\tau_0}N^{-3/4}.
\label{eq:app-known-curvature}
\end{equation}
The absence of a first-batch loss here demonstrates why $m/N$ was not
claimed as an unavoidable delay cost.

\subsection{Proofs of Proposition 2 and Euclidean Recovery}
\label{app:confidence}

\begin{proof}[Proof of \cref{prop:confidence} and its consequences]
Fix a completed horizon on $\mathcal E_\delta$. For any unit vector $u$,
Cauchy--Schwarz gives
\begin{align}
|u^\top e_T|
&=|(\Vhat_T^{-1/2}u)^\top(\Vhat_T^{1/2}e_T)|\notag\\
&\le\norm{\Vhat_T^{-1/2}u}_2\norm{e_T}_{\Vhat_T}
\le\beta_N(\delta)\sqrt{u^\top\Vhat_T^{-1}u}.
\label{eq:app-directional-proof}
\end{align}
There was no direction-dependent probability event, so this holds
simultaneously for all unit directions, including any direction
selected from the data. The same score event is used in
\cref{thm:information}.

For a realized vector, $\sup_{\norm{u}_2=1}|u^\top e_T|=\norm{e_T}_2$.
If $e_T\ne0$, equality is attained at $u=e_T/\norm{e_T}_2$; if
$e_T=0$, it is immediate. Also
\begin{equation}
\sup_{\norm{u}_2=1}u^\top\Vhat_T^{-1}u
=\lambda_{\max}(\Vhat_T^{-1})
=1/\lambda_{\min}(\Vhat_T).
\label{eq:app-ellipsoid-largest-radius}
\end{equation}
Taking the supremum in \eqref{eq:app-directional-proof} gives
\begin{equation}
\norm{e_T}_2\le\frac{\beta_N(\delta)}
{\sqrt{\lambda_{\min}(\Vhat_T)}}.
\label{eq:app-euclidean-ellipsoid}
\end{equation}
When $\varepsilon_{N,m}(\delta)<\gamma^\star$, the first inequality
of \eqref{eq:information-rate} gives
$\lambda_{\min}(\Vhat_T)\ge
N(\gamma^\star-\varepsilon_{N,m}(\delta))>0$.
Substitution proves \eqref{eq:l2-confidence}. Consistency is already
established in \eqref{eq:app-probability-one} independently of this
rate simplification.
\end{proof}

\subsection{Proofs of Proposition 3 and the Adaptive Risk Lower Bound}
\label{app:lower}

\begin{proof}[Proof of \cref{prop:risk-lower}]
\textit{Step 1: identify the likelihood of the adaptive experiment.}
Write $\pi_t(d\mathcal B_t\mid\mathcal D_{t-1})$ for a policy's
conditional batch-selection kernel. Let $\zeta$ denote the reference
measure for a single GLM response (counting or Lebesgue measure as
appropriate). On a realized history, the joint law factors as
\begin{equation}
d\Prob_\theta^\pi
=\prod_{t=1}^T\left\{
 \pi_t(d\mathcal B_t\mid\mathcal D_{t-1})
 \prod_{j=1}^mp_\theta(Y_{t,j}\mid x_{t,j})\,d\zeta(Y_{t,j})
\right\}.
\label{eq:app-adaptive-factorization}
\end{equation}
The kernels may be randomized but have no separate dependence on
$\theta$. Regularity and common response support therefore give the
likelihood ratio, almost surely under $\theta^\star$,
\begin{align}
\frac{d\Prob_\theta^\pi}{d\Prob_{\theta^\star}^\pi}
&=\prod_{i=1}^N\frac{p_\theta(Y_i\mid x_i)}
{p_{\theta^\star}(Y_i\mid x_i)}\notag\\
&=\exp\!\left\{\sum_{i=1}^N
\big[Y_ix_i^\top(\theta-\theta^\star)
-A(x_i^\top\theta)+A(x_i^\top\theta^\star)\big]\right\}.
\label{eq:app-adaptive-ratio-general}
\end{align}
For a fixed realized history, the policy factors cancel. The data
history is not differentiated through the query-selection algorithm.
Differentiating the log ratio at $\theta^\star$ gives $U_N$.
By \eqref{eq:app-score-cov}, the Fisher information of the entire
experiment is
\begin{equation}
\mathcal J_T^\pi=\E[U_NU_N^\top]
=\E\!\left[\sum_{i=1}^NG_{\theta^\star}(x_i)\right].
\label{eq:app-adaptive-fisher}
\end{equation}
The ridge term is absent because it is not part of the response
likelihood.

\textit{Step 2: bound expected information using a feasible measure.}
For a Borel subset $E\subseteq\mathcal X$, define the expected
empirical measure
\begin{equation}
\overline\nu_N(E)=\frac1N\sum_{i=1}^N\Prob_{\theta^\star}(x_i\in E).
\label{eq:app-expected-measure}
\end{equation}
Each summand is a probability measure, so $\overline\nu_N$ is one.
Linearity of integration yields
\begin{equation}
M(\overline\nu_N)
=\frac1N\sum_{i=1}^N\E[G_{\theta^\star}(x_i)]
=\frac{\mathcal J_T^\pi}{N}.
\label{eq:app-expected-feasible-information}
\end{equation}
Since $\overline\nu_N$ is feasible,
\begin{equation}
\lambda_{\min}(\mathcal J_T^\pi)
=N\lambda_{\min}(M(\overline\nu_N))\le N\gamma^\star.
\label{eq:app-expected-information-ceiling}
\end{equation}
No equality between $\lambda_{\min}(\E Z_N)$ and
$\E\lambda_{\min}(Z_N)$ is used.

\textit{Step 3: state precisely the estimator differentiation condition.}
For the competing estimator, define
$\widetilde e_T=\widetilde\theta_T-\theta^\star$ and
$\widetilde Q_T=\E[\widetilde e_T\widetilde e_T^\top]$.
The likelihood-differentiation condition in the main text is the
finite-horizon identity
\begin{equation}
D_T=\nabla_\theta\E_\theta[\widetilde\theta_T]\big|_{\theta=\theta^\star}
=\E_{\theta^\star}[\widetilde\theta_TU_N^\top]
=\E_{\theta^\star}[\widetilde e_TU_N^\top].
\label{eq:app-mean-score-identity}
\end{equation}
The last equality uses $\E U_N=0$. For arbitrary competing estimators,
this interchange of differentiation and integration is an explicit
condition, not a consequence of finite second moments alone.
For the proposed bounded ridge MLE it will be proved in
\cref{app:mean-derivative}.

\textit{Step 4: use a scalar information inequality without inversion.}
Choose a unit minimum eigenvector $v$ of $\mathcal J_T^\pi$.
If $D_Tv=0$, then $s_{\min}(D_T)=0$ and the claimed lower bound is
zero. Otherwise let $u=D_Tv/\norm{D_Tv}_2$. From
\eqref{eq:app-mean-score-identity},
\begin{align}
\norm{D_Tv}_2
&=u^\top D_Tv
=\E[(u^\top\widetilde e_T)(v^\top U_N)],\notag\\
\norm{D_Tv}_2^2
&\le\E[(u^\top\widetilde e_T)^2]\E[(v^\top U_N)^2]\notag\\
&=(u^\top\widetilde Q_Tu)(v^\top\mathcal J_T^\pi v)\notag\\
&\le\lambda_{\max}(\widetilde Q_T)
\lambda_{\min}(\mathcal J_T^\pi)
\le\lambda_{\max}(\widetilde Q_T)N\gamma^\star.
\label{eq:app-cr-direction}
\end{align}
Since $\norm{D_Tv}_2\ge s_{\min}(D_T)$ and
$\lambda_{\max}(\widetilde Q_T)
=\mathfrak R_{\mathrm{dir},T}(\widetilde\theta_T)^2$,
\eqref{eq:risk-lower-general} follows. This argument includes singular
$\mathcal J_T^\pi$: its inverse is never taken. In particular, a zero
information eigenvalue forces $D_Tv=0$ through the same inequality.

\textit{Step 5: derive the stated comparison classes.}
Local unbiasedness supplies $D_T=I$, so $s_{\min}(D_T)=1$ and
\eqref{eq:risk-lower} follows. If $D_T\to I$, then for every unit $v$,
\begin{equation}
1-\norm{D_T-I}_{\op}\le\norm{D_Tv}_2
\le1+\norm{D_T-I}_{\op}.
\label{eq:app-smallest-singular-limit}
\end{equation}
Taking the infimum over $v$ gives $s_{\min}(D_T)\to1$. Multiplication
of \eqref{eq:risk-lower-general} by $N\gamma^\star=T\gstar$ and
taking the lower limit proves \eqref{eq:asymptotic-lower}.
No finite-sample unbiasedness is imposed on the proposed estimator.
\end{proof}

\subsection{Proofs of Proposition 4 and Sharp Risk Attainment}
\label{app:risk}

In this subsection alone, $m$ is fixed, $\theta^\star$ is interior,
and the optimal information matrix $M^\star$ is unique as specified
in \cref{sec:sharp-risk}. Since $\lambda_{\min}(M^\star)=\gamma^\star>0$,
its inverse is well-defined. All counts below are completed counts
$N=mT$. The proof establishes the matrix MSE limit, its risk
consequences, and the mean-Jacobian claim after the proposition.

\subsubsection{Stabilization of information and the integrated Hessian}

Work first on the probability-one event from
\eqref{eq:app-probability-one}. The unregularized averages
$(Z_N-\rho I)/N$ all belong to the compact set $\mathcal C$ of
\eqref{eq:app-feasible-matrices}. Their minimum eigenvalues converge
to $\gamma^\star$. If a subsequence of these matrices converges to
$M_\infty$, continuity of the minimum eigenvalue gives
\begin{equation}
M_\infty\in\mathcal C,\qquad
\lambda_{\min}(M_\infty)=
\lim\lambda_{\min}((Z_N-\rho I)/N)=\gamma^\star.
\label{eq:app-optimal-cluster-point}
\end{equation}
By matrix uniqueness, $M_\infty=M^\star$.
To see that the full sequence converges, suppose it did not: there
would be an $\epsilon>0$ and a subsequence at distance at least
$\epsilon$ from $M^\star$. Compactness would supply a convergent
further subsequence, whose limit would be both optimal and at that
distance, a contradiction. Hence
\begin{equation}
Z_N/N\longrightarrow M^\star\qquad\text{a.s.}
\label{eq:app-matrix-stabilization}
\end{equation}
The vanishing ridge $\rho I/N$ does not change the limit.
This cluster-point argument is the use of uniqueness; it is not
assumed in the information-growth theorem.

For the integrated Hessian in \eqref{eq:app-integrated-hessian},
\begin{align}
\norm{K_N-Z_N}_{\op}
&\le\sum_{i=1}^N\int_0^1
 |w_{\theta^\star+s e_T}(x_i)-w_{\theta^\star}(x_i)|
 \norm{x_i}_2^2\,ds\notag\\
&\le\sum_{i=1}^N\int_0^1L_A s|x_i^\top e_T|\norm{x_i}_2^2\,ds
\le\frac12NC_F\norm{e_T}_2.
\label{eq:app-hessian-error}
\end{align}
Consistency and \eqref{eq:app-initial-tracking} therefore imply
\begin{equation}
K_N/N\to M^\star,
\qquad\Vhat_T/N\to M^\star
\qquad\text{a.s.}
\label{eq:app-hessian-stabilization}
\end{equation}
Choose $r_\theta>0$ so that the open ball with center $\theta^\star$
and radius $r_\theta$ is contained in $\operatorname{int}\Theta$.
Define the interior-estimate event
\begin{equation}
\mathcal I_N=\{\widehat\theta_T\in\operatorname{int}\Theta\}.
\label{eq:app-interior-event}
\end{equation}
Because $\norm{e_T}_2<r_\theta$ implies $\mathcal I_N$, consistency
also shows that this event occurs eventually on almost every path.

\subsubsection{Fourth moment of the normalized score}

For an arbitrary deterministic vector $v\in\R^d$, repeated
conditioning in \eqref{eq:app-predictable-noise} gives, for $a\in\R$,
\begin{align}
\E[e^{a v^\top U_n}]
&=\E\!\left[e^{a v^\top U_{n-1}}
 \E[e^{a(v^\top x_n)\xi_n}\mid\mathscr F_{n-1}]\right]\notag\\
&\le\E\!\left[e^{a v^\top U_{n-1}}
 e^{a^2\sigma^2(v^\top x_n)^2/2}\right]\notag\\
&\le e^{a^2\sigma^2L^2\norm{v}_2^2/2}
 \E[e^{a v^\top U_{n-1}}]\notag\\
&\le e^{a^2\sigma^2nL^2\norm{v}_2^2/2}.
\label{eq:app-score-mgf}
\end{align}
Let $s_N=U_N/\sqrt N$ denote the normalized score. For coordinate
$r$, exponential Markov inequality, optimized over $a>0$, gives
\begin{equation}
\Prob(|(U_N)_r|\ge u)
\le2\exp\!\left(-\frac{u^2}{2\sigma^2NL^2}\right)
\qquad(u>0).
\label{eq:app-coordinate-tail}
\end{equation}
For a nonnegative random variable, the tail integral formula follows
from integrating its pointwise indicator and applying Tonelli's
theorem. With $v_N=\sigma^2NL^2$, it yields
\begin{align}
\E[(U_N)_r^4]
&=\int_0^\infty4u^3\Prob(|(U_N)_r|>u)\,du\notag\\
&\le8\int_0^\infty u^3e^{-u^2/(2v_N)}\,du
=16v_N^2\int_0^\infty z e^{-z}\,dz
=16v_N^2.
\label{eq:app-coordinate-fourth}
\end{align}
The substitution is $z=u^2/(2v_N)$, for which
$u^3du=2v_N^2z\,dz$.
Since $(\sum_ra_r^2)^2\le d\sum_ra_r^4$,
\begin{equation}
\E\norm{s_N}_2^4
\le\frac d{N^2}\sum_{r=1}^d\E[(U_N)_r^4]
\le C_4,\qquad C_4=16d^2\sigma^4L^4.
\label{eq:app-score-fourth}
\end{equation}
In particular $\E\norm{s_N}_2^2\le C_4^{1/2}$ by Cauchy--Schwarz.

\subsubsection{Fourth moment of the normalized estimation error}

Define the well-conditioned event and a positive Hessian constant by
\begin{equation}
\mathcal G_N=\{\lambda_{\min}(Z_N)\ge N\gamma^\star/2\},
\qquad c_K=\frac{a_0\gamma^\star}{2b_0}.
\label{eq:app-good-event}
\end{equation}
For $N\ge2$, apply \cref{thm:information} with
$\delta_N=N^{-6}$. Since $m$ is fixed,
\begin{equation}
\varepsilon_{N,m}(N^{-6})
=C\left[\frac{1+7\log N}{N^{1/4}}+\frac mN\right]\to0.
\label{eq:app-polynomial-confidence}
\end{equation}
There is therefore a deterministic threshold $N_0=N_0(m)\ge2$ such
that, for every completed $N\ge N_0$,
\begin{equation}
\Prob(\mathcal G_N^c)\le N^{-6}.
\label{eq:app-rare-information}
\end{equation}
Indeed the theorem's event at that $N$ implies
\eqref{eq:app-refined-information}. The varying choice $\delta_N$
is used only for this marginal probability estimate; no claim that
all these confidence events hold together is required here.

From $K_N\succeq a_0H_N$ and $Z_N\preceq b_0H_N$,
\begin{equation}
K_N\succeq(a_0/b_0)Z_N\succeq c_KNI
\qquad\text{on }\mathcal G_N.
\label{eq:app-good-hessian}
\end{equation}
Combining this with the constrained inequality
\eqref{eq:app-variational} yields
\begin{align}
c_KN\norm{e_T}_2^2
&\le e_T^\top K_Ne_T
\le e_T^\top(U_N-\rho\theta^\star)
\le\norm{e_T}_2\norm{U_N-\rho\theta^\star}_2,\notag\\
\norm{\sqrt N e_T}_2
&\le c_K^{-1}\norm{s_N-\rho\theta^\star/\sqrt N}_2
\qquad\text{on }\mathcal G_N.
\label{eq:app-euclidean-score-bound}
\end{align}
As before, the zero-error case needs no division. This estimate is
valid even when the minimizer is on the boundary.

Let $z_N=\sqrt N e_T$ denote the scaled error. Using
$(a+b)^4\le8(a^4+b^4)$ and \eqref{eq:app-score-fourth},
\begin{equation}
\E[\norm{z_N}_2^4\mathbf1_{\mathcal G_N}]
\le8c_K^{-4}\left(C_4+\frac{\rho^4S^4}{N^2}\right).
\label{eq:app-good-estimator-fourth}
\end{equation}
On $\mathcal G_N^c$, compactness of $\Theta$ gives
$\norm{z_N}_2^4\le16S^4N^2$, so for $N\ge N_0$,
\begin{equation}
\E[\norm{z_N}_2^4\mathbf1_{\mathcal G_N^c}]
\le16S^4N^2\Prob(\mathcal G_N^c)
\le16S^4N^{-4}.
\label{eq:app-bad-estimator-fourth}
\end{equation}
For the finitely many completed $N<N_0$, use the bound
$16S^4N_0^2$. For example, the finite constant
\begin{equation}
C_e=\max\left\{
8c_K^{-4}(C_4+\rho^4S^4)+16S^4,\;16S^4N_0^2\right\}
\label{eq:app-estimator-moment-constant}
\end{equation}
satisfies
\begin{equation}
\sup_{T\ge1}\E\norm{\sqrt{mT}e_T}_2^4\le C_e.
\label{eq:app-estimator-fourth}
\end{equation}
This is a proved bound, not an assumption. In particular, for $K>0$,
\begin{equation}
\sup_T\E\!\left[\norm{z_N}_2^2
\mathbf1_{\{\norm{z_N}_2^2>K\}}\right]
\le\frac1K\sup_T\E\norm{z_N}_2^4\le C_e/K\to0,
\label{eq:app-uniform-integrability}
\end{equation}
which explicitly verifies uniform integrability of the squared
scaled errors.

\subsubsection{Mean-square linearization including exceptional events}

On $\mathcal I_N$ the constrained minimizer is interior, so its gradient
vanishes. Equation \eqref{eq:app-gradient-expansion} then becomes
\begin{equation}
K_Ne_T=U_N-\rho\theta^\star
\qquad\text{on }\mathcal I_N.
\label{eq:app-interior-score}
\end{equation}
Define the normalized Hessian $\overline K_N=K_N/N$, the intersection
$\mathcal A_N=\mathcal G_N\cap\mathcal I_N$, and the linearization
remainder
\begin{equation}
r_N=z_N-(M^\star)^{-1}s_N.
\label{eq:app-remainder-definition}
\end{equation}
Our aim is $\E\norm{r_N}_2^2\to0$, not merely convergence in
probability.

On $\mathcal A_N$, \eqref{eq:app-interior-score} gives
\begin{equation}
r_N=[\overline K_N^{-1}-(M^\star)^{-1}]s_N
-\overline K_N^{-1}\frac{\rho\theta^\star}{\sqrt N}.
\label{eq:app-remainder-on-good}
\end{equation}
The inverse difference has the exact identity
\begin{equation}
\overline K_N^{-1}-(M^\star)^{-1}
=\overline K_N^{-1}(M^\star-\overline K_N)(M^\star)^{-1}.
\label{eq:app-inverse-difference}
\end{equation}
On $\mathcal G_N$, \eqref{eq:app-good-hessian} gives
$\norm{\overline K_N^{-1}}_{\op}\le c_K^{-1}$; also
$\norm{(M^\star)^{-1}}_{\op}=1/\gamma^\star$.
Define the truncated difference
$\Delta_N=\mathbf1_{\mathcal A_N}
[\overline K_N^{-1}-(M^\star)^{-1}]$.
Then
\begin{equation}
\norm{\Delta_N}_{\op}\le c_K^{-1}+1/\gamma^\star,
\qquad
\Delta_N\to0\quad\text{a.s.}
\label{eq:app-truncated-inverse}
\end{equation}
The convergence follows from \eqref{eq:app-hessian-stabilization} and
\eqref{eq:app-inverse-difference}; the indicator cannot enlarge a
vanishing inverse difference. The displayed deterministic bound
justifies dominated convergence, giving
$\E\norm{\Delta_N}_{\op}^4\to0$.
Cauchy--Schwarz applied to the product in
\eqref{eq:app-remainder-on-good} yields
\begin{align}
\E[\norm{r_N}_2^2\mathbf1_{\mathcal A_N}]
&\le2\E[\norm{\Delta_N}_{\op}^2\norm{s_N}_2^2]
+\frac{2\rho^2S^2}{c_K^2N}\notag\\
&\le2(\E\norm{\Delta_N}_{\op}^4)^{1/2}C_4^{1/2}
+\frac{2\rho^2S^2}{c_K^2N}\longrightarrow0.
\label{eq:app-good-remainder-l2}
\end{align}

It remains to control the whole complement, including boundary
estimates. From \eqref{eq:app-remainder-definition},
\begin{equation}
\E\norm{r_N}_2^4
\le8\big(\E\norm{z_N}_2^4+
 (\gamma^\star)^{-4}\E\norm{s_N}_2^4\big)
\le C_r,\qquad C_r=8(C_e+(\gamma^\star)^{-4}C_4).
\label{eq:app-remainder-fourth}
\end{equation}
The chosen interior radius and Markov's inequality imply
\begin{equation}
\Prob(\mathcal I_N^c)
\le\Prob(\norm{e_T}_2\ge r_\theta)
\le\frac{\E\norm{z_N}_2^4}{N^2r_\theta^4}
\le\frac{C_e}{N^2r_\theta^4}.
\label{eq:app-boundary-probability}
\end{equation}
For $N\ge N_0$, combine this with
\eqref{eq:app-rare-information} and Cauchy--Schwarz:
\begin{align}
\E[\norm{r_N}_2^2\mathbf1_{\mathcal A_N^c}]
&\le(\E\norm{r_N}_2^4)^{1/2}
\Prob(\mathcal A_N^c)^{1/2}\notag\\
&\le C_r^{1/2}\left(N^{-6}+
\frac{C_e}{N^2r_\theta^4}\right)^{1/2}\longrightarrow0.
\label{eq:app-bad-remainder-l2}
\end{align}
Adding \eqref{eq:app-good-remainder-l2} and
\eqref{eq:app-bad-remainder-l2} proves
\begin{equation}
\E\left\|\sqrt N e_T-(M^\star)^{-1}\frac{U_N}{\sqrt N}\right\|_2^2
=\E\norm{r_N}_2^2\longrightarrow0.
\label{eq:app-l2-linearization}
\end{equation}
No central limit theorem or unproved moment interchange appears in
this argument.

\subsubsection{Matrix MSE and the two risk constants}

Since $\norm{(Z_N-\rho I)/N}_{\op}\le B$, matrix stabilization and
bounded convergence give
\begin{equation}
\E[s_Ns_N^\top]
=\E[(Z_N-\rho I)/N]\longrightarrow M^\star.
\label{eq:app-score-second-limit}
\end{equation}
More explicitly,
$\|\E[(Z_N-\rho I)/N]-M^\star\|_{\op}$ is at most
$\E\|(Z_N-\rho I)/N-M^\star\|_{\op}$, which tends to zero by
dominated convergence with upper bound $2B$.

Set $b_N^\star=(M^\star)^{-1}s_N$ for this paragraph.
Then $z_N=b_N^\star+r_N$, and expansion gives
\begin{align}
NQ_T
=\E[z_Nz_N^\top]
={}&(M^\star)^{-1}\E[s_Ns_N^\top](M^\star)^{-1}\notag\\
&+\E[b_N^\star r_N^\top]
+\E[r_N(b_N^\star)^\top]+\E[r_Nr_N^\top].
\label{eq:app-mse-expansion}
\end{align}
The last three terms vanish in operator norm. Indeed,
\begin{align}
\norm{\E[r_Nr_N^\top]}_{\op}
&\le\E\norm{r_N}_2^2\to0,\notag\\
\norm{\E[b_N^\star r_N^\top]}_{\op}
&\le\E[\norm{b_N^\star}_2\norm{r_N}_2]\notag\\
&\le\frac1{\gamma^\star}
(\E\norm{s_N}_2^2)^{1/2}(\E\norm{r_N}_2^2)^{1/2}
\le\frac{C_4^{1/4}}{\gamma^\star}
(\E\norm{r_N}_2^2)^{1/2}\to0.
\label{eq:app-mse-cross-terms}
\end{align}
The transposed cross term has the same norm. Thus
\begin{equation}
NQ_T\longrightarrow(M^\star)^{-1}M^\star(M^\star)^{-1}
=(M^\star)^{-1}.
\label{eq:app-mse-limit}
\end{equation}
This proves the formal statement of \cref{prop:efficiency}.
Because $N=mT$ and $J^\star=mM^\star$,
\begin{equation}
TQ_T=(NQ_T)/m\longrightarrow (M^\star)^{-1}/m=(J^\star)^{-1}.
\label{eq:app-per-round-mse}
\end{equation}
Continuity of the largest eigenvalue and trace, and the risk identities
in \eqref{eq:risk-matrix}, now yield
\begin{align}
T\gstar\,\Rdir^2
&=\gstar\lambda_{\max}(TQ_T)
\longrightarrow\gstar\lambda_{\max}((J^\star)^{-1})=1,\notag\\
T\Rtwo^2
&=\tr(TQ_T)\longrightarrow\tr((J^\star)^{-1}).
\label{eq:app-risk-constant-calculation}
\end{align}
For completeness, the continuity bounds are
$|\lambda_{\max}(A_1)-\lambda_{\max}(A_2)|
\le\norm{A_1-A_2}_{\op}$ and
$|\tr(A_1-A_2)|\le d\norm{A_1-A_2}_{\op}$.
If $\lambda_1,\ldots,\lambda_d$ are the eigenvalues of $J^\star$,
then $\min_r\lambda_r=\gstar$, so
\begin{equation}
\frac1{\gstar}\le\sum_{r=1}^d\frac1{\lambda_r}
=\tr((J^\star)^{-1})\le\frac d{\gstar}.
\label{eq:app-trace-risk-range}
\end{equation}
These computations identify the Euclidean constant but do not say
that an E-optimal design minimizes the trace among all designs.

\subsubsection{Differentiating the proposed estimator's mean}
\label{app:mean-derivative}

We now prove the remaining comparison-class assertion for the ridge
MLE. Fix a finite completed horizon $N$. The policy and estimator are
held fixed as the statistical parameter varies. For a vector
$h\in\R^d$ with $\norm{h}_2<r_\theta$, define the likelihood ratio
$L_N(h)=d\Prob_{\theta^\star+h}/d\Prob_{\theta^\star}$.
The factorization \eqref{eq:app-adaptive-ratio-general} gives
\begin{equation}
L_N(h)=\exp\{h^\top U_N-R_N(h)\},
\label{eq:app-likelihood-ratio}
\end{equation}
where the scalar Taylor remainder is
\begin{equation}
R_N(h)=\sum_{i=1}^N
\left[A(x_i^\top\theta^\star+x_i^\top h)
-A(x_i^\top\theta^\star)-A'(x_i^\top\theta^\star)x_i^\top h\right].
\label{eq:app-likelihood-remainder}
\end{equation}
Taylor's formula in integral form yields
\begin{align}
R_N(h)
&=\sum_{i=1}^N\int_0^1(1-s)
 w_{\theta^\star+sh}(x_i)(x_i^\top h)^2\,ds,\notag\\
0\le R_N(h)&\le\tfrac12NB\norm{h}_2^2,\notag\\
\nabla_h R_N(h)
&=\sum_{i=1}^N\int_0^1
 w_{\theta^\star+sh}(x_i)x_ix_i^\top h\,ds,\notag\\
\norm{\nabla_hR_N(h)}_2&\le NB\norm{h}_2.
\label{eq:app-likelihood-remainder-bounds}
\end{align}
All intermediate parameters belong to $\Theta$ by the choice of
radius. Therefore, for a fixed $r_0\in(0,r_\theta)$ and
$\norm{h}_2\le r_0$,
\begin{equation}
\norm{\nabla_hL_N(h)}_2
=L_N(h)\norm{U_N-\nabla_hR_N(h)}_2
\le e^{r_0\norm{U_N}_2}(\norm{U_N}_2+NBr_0).
\label{eq:app-derivative-envelope}
\end{equation}
We verify that this envelope is integrable rather than assume it.
For any $a>0$, the scalar inequalities
$\norm{U_N}_2\le\sum_r|(U_N)_r|$ and
$e^{a|x|}\le e^{ax}+e^{-ax}$ imply
\begin{align}
\E e^{a\norm{U_N}_2}
&\le\sum_{v\in\{-1,1\}^d}\E e^{a v^\top U_N}
\le2^d\exp(a^2\sigma^2NL^2d/2)<\infty.
\label{eq:app-vector-exponential-envelope}
\end{align}
Here \eqref{eq:app-score-mgf} applies since $\norm{v}_2^2=d$.
For $u\ge0$, $u\le e^u$, so
\begin{equation}
e^{r_0u}(u+NBr_0)
\le(1+NBr_0)e^{(r_0+1)u}.
\label{eq:app-envelope-integrable}
\end{equation}
Equations \eqref{eq:app-vector-exponential-envelope} and
\eqref{eq:app-envelope-integrable} prove integrability of
\eqref{eq:app-derivative-envelope} at this fixed horizon.

At $h=0$, $L_N(0)=1$ and $\nabla_hL_N(0)=U_N$.
The fundamental theorem of calculus along the line segment gives
\begin{equation}
\frac{|L_N(h)-1-h^\top U_N|}{\norm{h}_2}
\le\int_0^1\norm{\nabla_hL_N(sh)-U_N}_2\,ds
\longrightarrow0\quad(h\to0)
\label{eq:app-frechet-likelihood}
\end{equation}
on each path. The right side is bounded by the integrable envelope in
\eqref{eq:app-derivative-envelope} plus $\norm{U_N}_2$.
Because $\norm{\widehat\theta_T}_2\le S$, dominated convergence
therefore shows the vector-valued Fr\'echet expansion
\begin{align}
a_T(\theta^\star+h)
&=\E_{\theta^\star}[\widehat\theta_T L_N(h)]\notag\\
&=a_T(\theta^\star)
 +\E_{\theta^\star}[\widehat\theta_TU_N^\top]h
 +o(\norm{h}_2).
\label{eq:app-mean-differentiation-expansion}
\end{align}
Thus the mean Jacobian exists and satisfies
\begin{equation}
D_T=\E[\widehat\theta_TU_N^\top]=\E[e_TU_N^\top],
\label{eq:app-mle-mean-derivative}
\end{equation}
where the second equality uses $\E U_N=0$.
This establishes the differentiation identity for each finite $T$
before any asymptotic limit is taken.

Finally, substitute $z_N=(M^\star)^{-1}s_N+r_N$ into this identity:
\begin{align}
D_T
&=\E[z_Ns_N^\top]
=(M^\star)^{-1}\E[s_Ns_N^\top]+\E[r_Ns_N^\top],\notag\\
\norm{\E[r_Ns_N^\top]}_{\op}
&\le(\E\norm{r_N}_2^2)^{1/2}
 (\E\norm{s_N}_2^2)^{1/2}\longrightarrow0.
\label{eq:app-derivative-remainder}
\end{align}
Using \eqref{eq:app-score-second-limit},
\begin{equation}
D_T\longrightarrow(M^\star)^{-1}M^\star=I.
\label{eq:app-derivative-limit}
\end{equation}
The same linearization also gives
\begin{equation}
\norm{\sqrt N\,\E e_T}_2
=\norm{\E r_N}_2
\le(\E\norm{r_N}_2^2)^{1/2}\longrightarrow0,
\label{eq:app-asymptotic-bias}
\end{equation}
since $\E s_N=0$. Finite-sample unbiasedness is not asserted.
Together with \cref{prop:risk-lower}, the established mean-Jacobian
limit and the sharp directional limit prove attainment within the
stated comparison class. This completes all claims accompanying
\cref{prop:efficiency}.

\subsection{Details of the Alternative Temperature Schedule}
\label{app:temperature}

The preceding results use $\tau_t=\tau_0(mt)^{3/4}$. The alternative
$\tau_t=\tau_0t^{3/4}$ has the same fixed-$m$ schedule family after
rescaling its free coefficient, but is a different numerical policy
when $\tau_0$ is held fixed across batch sizes. We give the full rate
calculation for this alternative rather than identifying the two
schedules without qualification.

Because $\tau_0t^{3/4}\le\tau_0(mt)^{3/4}$, the excitation proof still
holds with the same $c_E,N_E$. The confidence and preliminary
localization proofs do not otherwise use the temperature, so the
same $C_L$ applies. For the round-scaled schedule,
\begin{align}
\tau_T&=\tau_0T^{3/4}=\tau_0m^{-3/4}N^{3/4},\notag\\
m\sum_{t=1}^T\frac1{\tau_t}
&\le\frac{4m}{\tau_0}T^{1/4}
=\frac4{\tau_0}m^{3/4}N^{1/4},\notag\\
\frac{n_{t,j}}{\tau_t}
&\le\frac{mt}{\tau_0t^{3/4}}
=\tau_0^{-1}m^{3/4}(mt)^{1/4}.
\label{eq:app-round-scaled-components}
\end{align}
For $t\ge2$, the unchanged localization estimate now gives
\begin{equation}
\ell_{t,j}\le C_s\Lambda_N
\left[(mt)^{-1/2}+\frac{2B}{\tau_0}m^{3/4}(mt)^{-1/4}\right].
\label{eq:app-round-scaled-loss}
\end{equation}
Summing exactly as in \eqref{eq:app-power-sums} yields
\begin{equation}
\sum_{t=1}^T\sum_{j=1}^m\ell_{t,j}
\le Bm+2C_s\Lambda_NN^{1/2}
+\frac{8BC_s}{3\tau_0}m^{3/4}\Lambda_NN^{3/4}.
\label{eq:app-round-scaled-loss-sum}
\end{equation}
Substitution into the same master inequality
\eqref{eq:app-master-gap}, followed by division by $N$, gives
\begin{align}
0\le\frac{\rho+N\gamma^\star-\lambda_{\min}(Z_N)}N
\le{}&\tau_0\log d\,m^{-3/4}N^{-1/4}
+\frac{4B^2}{\tau_0}m^{3/4}N^{-3/4}+B\frac mN\notag\\
&+2C_s\Lambda_NN^{-1/2}
+\frac{8BC_s}{3\tau_0}m^{3/4}\Lambda_NN^{-1/4}.
\label{eq:app-round-scaled-decomposition}
\end{align}
Terminal plug-in error is still at most
$C_FC_L\Lambda_N/\sqrt N$ after normalization, by
\eqref{eq:app-initial-tracking}. Since $m\ge1$, $N\ge1$, and
$\Lambda_N\ge1$, the constant in \eqref{eq:app-main-constant}
(or a larger constant independent of $m$) bounds both rate errors by
\begin{equation}
\varepsilon^{\mathrm{round}}_{N,m}(\delta)
=C\left[m^{3/4}\frac{\log(eN/\delta)}{N^{1/4}}+\frac mN\right].
\label{eq:app-old-schedule-bound}
\end{equation}
For fixed $m$, this bound vanishes. Once it is at most
$\gamma^\star/2$, the refined localization calculation
\eqref{eq:app-refined-localization} applies unchanged. The
probability-one argument, matrix stabilization, fourth moments,
mean-square linearization, and mean differentiation then proceed
with this alternative rate error; for the fourth-moment proof,
$\varepsilon^{\mathrm{round}}_{N,m}(N^{-6})\to0$ for fixed $m$.
Thus the same fixed-$m$ asymptotic conclusions are retained.

Neither temperature analysis proves the zero-temperature hard policy,
an approximate finite-candidate oracle, a finite pool merely contained
in the ball, misspecified responses, or general Poisson noise under
the stated sub-Gaussian condition. These are not being introduced as
silent extensions of the theorem's scope.

\subsection{Range-Based Finite-Budget Temperature Calibration}
\label{app:range-calibration}

This subsection derives the constant $1/8$ in the range-based
reference temperature and explains the prescribed-budget calibration
in \eqref{eq:temperature}. The resulting optimization statement
concerns a spectral-smoothing penalty within a specified schedule
family, not the estimation error of the resulting adaptive policy.

\subsubsection{A range-dependent potential increment}

For this derivation, let $R>0$ bound the spectrum of each information
increment: an increment $D$ satisfies $0\preceq D\preceq RI$.
We use the Gibbs matrix $P_\tau$ and potential $\Phi_\tau$ defined
in \eqref{eq:gibbs} and \eqref{eq:potential}.

\begin{lemma}[Range-dependent potential increment]
\label{lem:range-potential-increment}
For every symmetric matrix $A$, every $D$ satisfying
$0\preceq D\preceq RI$, and every $\tau>0$,
\begin{equation}
\Phi_\tau(A+D)-\Phi_\tau(A)
\ge
\tr\!\left(P_\tau(A)D\right)-\frac{R^2}{8\tau}.
\label{eq:range-potential-increment}
\end{equation}
No commutation between $A$ and $D$ is required.
\end{lemma}

\begin{proof}
Write $z_A=\tr e^{-A/\tau}$ and $P=P_\tau(A)$.
The Golden--Thompson inequality \citep{Golden1965,Thompson1965}
gives
\begin{equation}
\frac{\tr e^{-(A+D)/\tau}}{z_A}
\le
\frac{\tr\!\left(e^{-A/\tau}e^{-D/\tau}\right)}{z_A}
=
\tr\!\left(Pe^{-D/\tau}\right).
\label{eq:range-golden-thompson}
\end{equation}
Let $r_1,\ldots,r_d\in[0,R]$ be the eigenvalues of $D$, with
orthonormal eigenvectors $v_1,\ldots,v_d$. Define
$q_k=v_k^\top Pv_k$. These numbers are nonnegative and sum to one.
Consequently,
\begin{equation}
\tr\!\left(Pe^{-D/\tau}\right)
=\sum_{k=1}^d q_k e^{-r_k/\tau},
\qquad
\tr(PD)=\sum_{k=1}^d q_k r_k.
\end{equation}
Hoeffding's lemma \citep{Hoeffding1963}, applied to the scalar
random variable taking value $r_k$ with probability $q_k$, yields
\begin{equation}
\log\!\left(\sum_{k=1}^d q_k e^{-r_k/\tau}\right)
\le
-\frac{\tr(PD)}{\tau}+\frac{R^2}{8\tau^2}.
\label{eq:range-hoeffding}
\end{equation}
Combining the last inequality with
\eqref{eq:range-golden-thompson} and multiplying by $-\tau$
proves \eqref{eq:range-potential-increment}.
\end{proof}

\subsubsection{The within-batch smoothing penalty}

Fix the fitted parameter during one batch. Let $W_0$ be its starting
plug-in information matrix, let $D_j$ be the information increment
of its $j$th selected query, and write $W_j=W_{j-1}+D_j$.
For a fixed batch temperature $\tau$, denote the selected query's
score by
\begin{equation}
s_j=\tr\!\left(P_\tau(W_{j-1})D_j\right),
\qquad j=1,\ldots,m.
\end{equation}
If $0\preceq D_j\preceq RI$ for every $j$, summing
Lemma~\ref{lem:range-potential-increment} gives
\begin{equation}
\Phi_\tau(W_m)-\Phi_\tau(W_0)
\ge
\sum_{j=1}^m s_j-\frac{mR^2}{8\tau}.
\label{eq:range-batch-potential}
\end{equation}
The soft-min bracket \eqref{eq:app-soft-min-bracket} also gives
\begin{equation}
\lambda_{\min}(W_m)-\lambda_{\min}(W_0)
\ge
\Phi_\tau(W_m)-\Phi_\tau(W_0)-\tau\log d.
\end{equation}
Define the resulting scalar penalty by
\begin{equation}
\Psi_m(\tau)=\tau\log d+\frac{mR^2}{8\tau}.
\label{eq:range-penalty}
\end{equation}
Then
\begin{equation}
\lambda_{\min}(W_m)-\lambda_{\min}(W_0)
\ge
\sum_{j=1}^m s_j-\Psi_m(\tau).
\label{eq:range-batch-eigenvalue}
\end{equation}
This is a pathwise statement for the frozen-parameter virtual batch.
It neither identifies $W_m$ with the information recomputed after
refitting nor requires the query matrices to commute.

For $d\ge2$, the derivatives of the penalty are
\begin{equation}
\Psi_m'(\tau)=\log d-\frac{mR^2}{8\tau^2},
\qquad
\Psi_m''(\tau)=\frac{mR^2}{4\tau^3}>0.
\end{equation}
Thus its unique unconstrained minimizer is
\begin{equation}
\tau_{\mathrm{range}}^{\circ}
:=\argmin_{\tau>0}\Psi_m(\tau)
=R\sqrt{\frac{m}{8\log d}}.
\label{eq:range-unfloored-minimizer}
\end{equation}
In particular, the factor $1/8$ follows from the range-dependent
exponential-moment bound, while the factor $m$ follows from summing
one increment bound per virtual update. The numerical reference
in \eqref{eq:temperature} is equivalently
\begin{equation}
\tau_{\mathrm{range}}
=\argmin_{\tau\ge10^{-8}}\Psi_m(\tau)
=\max\!\left\{10^{-8},\tau_{\mathrm{range}}^{\circ}\right\}.
\end{equation}
Minimizing this penalty alone does not maximize the full right-hand
side of \eqref{eq:range-batch-eigenvalue}: the selected queries and
their scores can also depend on $\tau$.

Taking $R=B$ from \eqref{eq:app-constants} gives the sharper
remainder $mB^2/(8\tau_t)$ in the true-state batch-gain argument of
Lemma~\ref{lem:batch}, with its existing score-discrepancy terms
unchanged. The preceding proofs retain their valid, looser remainder
$mB^2/\tau_t$; none of their displayed constants is being silently
replaced.

\subsubsection{Calibration within a capped power-law family}

The power $3/4$ retains the schedule shape used in the information-rate
analysis. To choose its coefficient over a prescribed budget, we impose
the design constraint that the temperature does not exceed the
range-based reference before that budget is exhausted. This cap is a
calibration choice, not a necessary condition of
Theorem~\ref{thm:information}.

Let $T_{\max}\ge1$ be the prescribed number of batches and
$N_{\max}=mT_{\max}$. For a positive coefficient $c$, consider
$\tau_t(c)=c(mt)^{3/4}$. The admissible coefficients satisfy
\begin{equation}
0<c\le\frac{\tau_{\mathrm{range}}}{N_{\max}^{3/4}},
\label{eq:range-coefficient-class}
\end{equation}
which is equivalent to $\tau_t(c)\le\tau_{\mathrm{range}}$
for every $1\le t\le T_{\max}$.

\begin{proposition}[Coefficient calibration within the capped family]
\label{prop:range-coefficient}
Suppose $d\ge2$, $R>0$, and the numerical floor is inactive, so
$\tau_{\mathrm{range}}=\tau_{\mathrm{range}}^{\circ}$.
Among the coefficients in \eqref{eq:range-coefficient-class},
\begin{equation}
c_\star=\frac{\tau_{\mathrm{range}}}{N_{\max}^{3/4}}
\end{equation}
uniquely minimizes $\Psi_m(\tau_t(c))$ for every
$t=1,\ldots,T_{\max}$. Its schedule is
\begin{equation}
\tau_t(c_\star)
=\tau_{\mathrm{range}}
\left(\frac{mt}{N_{\max}}\right)^{3/4}
=\tau_{\mathrm{range}}
\left(\frac{t}{T_{\max}}\right)^{3/4}.
\label{eq:range-anchored-schedule}
\end{equation}
\end{proposition}

\begin{proof}
The schedule is increasing in $t$, so the last batch supplies the
strongest cap and gives \eqref{eq:range-coefficient-class}.
The function $\Psi_m$ is strictly decreasing on
$(0,\tau_{\mathrm{range}}^{\circ}]$. For each fixed $t$, the
largest feasible coefficient gives the largest feasible temperature
in that interval, and therefore the smallest penalty. That
coefficient is $c_\star$.
\end{proof}

The proposition concerns only the scalar penalty in the stated
capped family. It does not establish optimality over unrestricted
temperature schedules, over adaptive query sequences, or for
Euclidean or worst-direction estimation risk.

\subsubsection{Choice of range and numerical conventions}

For the GLM information in \eqref{eq:fisher},
\begin{equation}
\norm{G_\theta(x)}_{\op}
=w_\theta(x)\norm{x}_2^2.
\end{equation}
Hence a valid uniform choice for the preceding range-bound derivation
is $R=\overline wL^2$ under the radius and curvature conditions in
Section~\ref{sec:conditions}. On unit-radius pools, one may take
$R=1$ for unit-variance Gaussian responses and $R=1/4$ for
Bernoulli-logistic responses.

For an unnormalized finite pool, let
$L_{\mathcal X}=\max_{x\in\mathcal X}\norm{x}_2$ be its actual
radius. The corresponding sufficient range bounds are
$R=L_{\mathcal X}^2$ and $R=L_{\mathcal X}^2/4$, respectively.
Using $R=1$ or $R=1/4$ without a unit-radius condition is instead a
family-reference calibration: the temperature remains well-defined,
but these values need not bound the information increments, so the
range-bound interpretation is not asserted for that implementation.
The uniform constant $B=\overline wL^2$ in the theoretical proofs
is unchanged.

The floor $10^{-8}$ is applied to the terminal reference temperature
only. No additional clipping of each $\tau_t$ is part of
\eqref{eq:range-anchored-schedule}. If the floor is active, the
schedule is still defined by \eqref{eq:temperature}, but
Proposition~\ref{prop:range-coefficient} is not invoked.
For $d=1$, $P_\tau(W)=1$ for every positive temperature, so the
spectral acquisition rule is temperature-independent and no
$\log d$-based calibration is needed. When a warm start is used,
$mt$ and $N_{\max}$ count adaptive responses after that warm start;
the warm-start information remains in the initial state.

\subsubsection{Relation to the fixed-coefficient asymptotic analysis}

Theorem~\ref{thm:information} analyzes a power-law schedule with a
fixed coefficient $\tau_0>0$. For each fixed choice of
$m$, $R$, and $N_{\max}$, \eqref{eq:temperature} supplies such a
coefficient. If this coefficient is retained and the same power law
is continued, the existing results apply under their original
full-ball, estimator, and response assumptions. The temperature
cap above applies only up to the prescribed calibration budget;
continuation does not clip the temperature at $\tau_{\mathrm{range}}$.

By contrast, increasing $N_{\max}$ and recalibrating $\tau_0$ for
every new run defines a family of different policies. The existing
bounds depend on $\tau_0$, including through $1/\tau_0$, and do
not by themselves establish uniform asymptotic optimality for that
horizon-dependent family. Likewise, recalibration across batch sizes
does not preserve a claim of constants uniform in $m$ without
tracking this dependence. The local calibration result does not
extend the full-ball guarantees to finite-pool sampling without
replacement or to misspecified response models.

\end{document}
